\documentclass[12pt,a4paper]{article}

% ===== Packages =====
\usepackage[utf8]{inputenc}
\usepackage[T1]{fontenc}
\usepackage{amsmath,amssymb,amsfonts}
\usepackage[margin=2.5cm]{geometry}
\usepackage[shortlabels]{enumitem}
\usepackage{hyperref}
\usepackage{xcolor}
\usepackage{booktabs}
\usepackage{longtable}

% ===== Custom commands =====
\newcommand{\Tr}{\operatorname{Tr}}
\newcommand{\Lam}{\Lambda}
\newcommand{\half}{\tfrac{1}{2}}
\newcommand{\task}[1]{\noindent\textbf{Task:} #1\par}
\newcommand{\method}[1]{\noindent\textbf{Method:} #1\par}
\newcommand{\outcome}[1]{\noindent\textbf{Expected outcome:} #1\par}
\newcommand{\deps}[1]{\noindent\textbf{Dependencies:} #1\par}

\definecolor{proven}{RGB}{0,100,0}
\definecolor{near}{RGB}{0,0,160}
\definecolor{mid}{RGB}{140,80,0}
\definecolor{long}{RGB}{160,0,0}

% ===== Title =====
\title{\textbf{Spectral Causal Theory (SCT):\\[4pt]
Research Roadmap and Task Prioritization}}
\author{David Alfyorov}
\date{March 2026}

\begin{document}
\maketitle
\begin{center}
\small\textit{Credits: Aliaksandr Samatyia contributed research-assistance and workflow support.}
\end{center}
\tableofcontents
\newpage

% ==========================================================================
\section{\textcolor{proven}{What Has Been Proven}}
\label{sec:proven}
% ==========================================================================

This section summarizes the results established in Part~VII of the theory development.
All results below are derived, not postulated.

\subsection{Classical limit: Einstein--Hilbert action from spectral action}

The spectral action principle
\begin{equation}\label{eq:spectral-action}
  S_{\text{spec}} = \Tr\!\left(f\!\left(\frac{D^2}{\Lam^2}\right)\right)
\end{equation}
where $D$ is the Dirac operator of the spectral triple $(\mathcal{A}, \mathcal{H}, D)$
and $f$ is a smooth cutoff function, reproduces in the IR limit the gravitational action:
\begin{equation}\label{eq:EH-limit}
  S_{\text{spec}} \;\xrightarrow{\;\text{IR}\;}\;
  \frac{1}{16\pi G}\int d^4x\,\sqrt{g}\,
  \bigl(-R + 2\Lam_{\text{cc}}\bigr)
  + c_1 \int d^4x\,\sqrt{g}\, C_{\mu\nu\rho\sigma}C^{\mu\nu\rho\sigma}
  + c_2 \int d^4x\,\sqrt{g}\, R^2 + \cdots
\end{equation}
with the cosmological constant $\Lam_{\text{cc}}$ and curvature-squared coefficients
$c_1$, $c_2$ determined by moments of $f$.

\textbf{Status:} Proven. The coefficients are fixed by the spectral triple,
not free parameters.

\subsection{Quantum potential and the coefficient $\beta$}

The coefficient in the one-loop effective potential:
\begin{equation}\label{eq:beta}
  \beta = \frac{41}{10\pi}
\end{equation}
is reproduced automatically because the IR limit of SCT matches GR.
This is a consistency check: SCT does not introduce new free parameters at one loop
but reproduces the known result from the spectral structure.

\textbf{Status:} Proven as consequence of IR matching.

\subsection{First unique prediction: the ratio $c_1/c_2$}

The ratio of curvature-squared coefficients:
\begin{equation}\label{eq:ratio}
  \frac{c_1}{c_2} = -\frac{1}{3}
\end{equation}
is \emph{fixed} by the spectral triple. In generic higher-derivative gravity,
$c_1$ and $c_2$ are independent parameters. SCT predicts a definite relation,
which is in principle testable via:
\begin{itemize}
  \item gravitational wave dispersion in the high-frequency regime,
  \item modifications to black hole quasinormal mode spectra,
  \item corrections to Newtonian potential at short distances.
\end{itemize}

\textbf{Status:} Proven for the minimal spectral triple.
\textbf{Caveat:} the sign correction $E = -R/4$ for Dirac fields
(established in NT-1) changes $\hat{P} = -5R/12$ and gives
$\beta_R^{\text{Dirac}} = 0$. The ratio $c_1/c_2$ for the full
SM spectrum (with $N_s = 4$, $N_f = 45$, $N_v = 12$) must be
re-derived in NT-1b, step~(e). The value $-1/3$ may shift.

\subsection{Nonlocal form factors: the master function}

All five universal form factors arising in the one-loop gravitational effective action
are generated by a single master function:
\begin{equation}\label{eq:master}
  f(\xi) = \int_0^1 d\alpha\; e^{-\alpha(1-\alpha)\,\xi}
\end{equation}
The five form factors $F_1,\ldots,F_5$ are expressible as combinations of $f(\xi)$
and its derivatives. The structure mirrors Barvinsky--Vilkovisky nonlocal heat kernel
technology but here arises from spectral data rather than being postulated.

\textbf{Status:} Proven. Mathematical infrastructure established.

% ==========================================================================
\section{\textcolor{near}{Near-Term Tasks}}
\label{sec:near}
% ==========================================================================

These tasks are the immediate priorities. Each builds directly on the
proven results of Section~\ref{sec:proven}.

\subsection{NT-1: Explicit computation of $F_1(\Box/\Lam^2)$ and $F_2(\Box/\Lam^2)$}
\label{task:NT1}

\task{Compute the exact (non-expanded) form factors $F_1(\Box/\Lam^2)$
and $F_2(\Box/\Lam^2)$ from the full spectral action~\eqref{eq:spectral-action},
without truncating to the asymptotic heat kernel expansion.}

\method{Two complementary approaches:
\begin{enumerate}[(a)]
  \item \textbf{Zeta-function method.} Express the spectral action via the
    spectral zeta function $\zeta_D(s) = \Tr(|D|^{-2s})$. The form factors
    are obtained from the analytic continuation:
    \begin{equation}\label{eq:zeta-method}
      F_k(z) = \frac{1}{\Gamma(s_k)} \int_0^\infty dt\; t^{s_k - 1}\,
      \tilde{f}_k(t)\, e^{-t\,z}, \quad z = \Box/\Lam^2
    \end{equation}
    where $\tilde{f}_k$ are determined by the cutoff function $f$ and the
    Seeley--DeWitt coefficients of $D^2$.

  \item \textbf{Nonlocal heat kernel.} Use the Barvinsky--Vilkovisky technique
    extended to the full (resummed) heat kernel:
    \begin{equation}\label{eq:heat-kernel}
      \Tr\!\left(f\!\left(\frac{D^2}{\Lam^2}\right)\right)
      = \int_0^\infty ds\;\hat{f}(s)\,\Tr\!\left(e^{-s D^2/\Lam^2}\right)
    \end{equation}
    with $\hat{f}$ the inverse Laplace transform of $f$. Resum the curvature
    expansion to obtain closed-form $F_1$, $F_2$.
\end{enumerate}
}

\outcome{Explicit analytic expressions for $F_1(z)$ and $F_2(z)$ valid for
all $z \geq 0$, including:
\begin{itemize}
  \item IR limit ($z \to 0$): must reproduce Eq.~\eqref{eq:EH-limit} with ratio~\eqref{eq:ratio}.
  \item UV behavior ($z \to \infty$): determines whether the theory is asymptotically free,
    safe, or finite.
  \item Analyticity domain in the complex $z$-plane.
\end{itemize}
}

\deps{None (builds on proven results). This is the highest-priority task.}

\textbf{Status (2026-03-11):} \textcolor{proven}{COMPLETE.} Form factors derived
for the massless Dirac operator. Seven rounds of independent audit (274+
quantitative checks, all PASS). Deep verification via systematic multi-method checks
with 95--98\% confidence. See \texttt{theory/derivations/NT1\_form\_factors.tex}.

\subsection{NT-1b: Form factors for scalar and vector fields}
\label{task:NT1b}

\task{Extend the NT-1 computation to scalar (minimal and conformal coupling)
and vector (gauge) fields, completing the full Standard Model spectrum.}

\method{
The CZ form factors~\cite{CZ2012} are universal for any Laplace-type operator
$\Delta = -(g^{\mu\nu}\nabla_\mu\nabla_\nu + E)$.
Only the endomorphism $E$ and connection curvature $\Omega_{\mu\nu}$ differ
between field types:
\begin{center}
\begin{tabular}{lccc}
\toprule
\textbf{Field} & $E$ & $\Omega_{\mu\nu}$ & \textbf{tr(Id)} \\
\midrule
Scalar ($\xi = 0$) & $0$ & $0$ & $1$ \\
Scalar ($\xi = 1/6$) & $-R/6$ & $0$ & $1$ \\
Dirac (massless) & $-R/4$ & $R_{\mu\nu\rho\sigma}\gamma^{\rho}\gamma^{\sigma}/4$ & $4$ \\
Vector (gauge) & $R_{\mu\nu}$ & $R_{\mu\nu\rho\sigma}$ & $d$ \\
\bottomrule
\end{tabular}
\end{center}
\begin{enumerate}[(a)]
  \item For each field type, compute the Weyl-basis form factors
    $h_C^{(s)}(x)$ and $h_R^{(s)}(x)$ using the CZ combination formulas
    with the appropriate traces.
  \item Verify local limits against known $\beta$-function coefficients:
    \begin{equation}
      \beta_W^{\text{scalar}} = \frac{1}{120}, \quad
      \beta_W^{\text{Dirac}} = \frac{1}{20}, \quad
      \beta_W^{\text{vector}} = \frac{1}{10}
    \end{equation}
  \item Construct the total form factors for the SM spectrum:
    \begin{equation}\label{eq:total-ff}
      F_i^{\text{total}}(z) = N_s\,F_i^{\text{scalar}}(z)
      + N_f\,F_i^{\text{Dirac}}(z) + N_v\,F_i^{\text{vector}}(z)
    \end{equation}
    with $N_s = 4$, $N_f = 45$, $N_v = 12$.
  \item Compute the total propagator and verify that the UV behavior of the
    combined form factors is consistent with the individual-field results.
  \item \textbf{Ratio $c_1/c_2$ re-derivation.} Compute the total local limits
    $c_i^{\text{total}} = N_s\,F_i^{\text{scalar}}(0) + N_f\,F_i^{\text{Dirac}}(0)
    + N_v\,F_i^{\text{vector}}(0)$ and determine the corrected ratio
    $c_1^{\text{total}}/c_2^{\text{total}}$.
    \textbf{Critical:} the sign correction $E = -R/4$ for Dirac (established in NT-1)
    gives $\hat{P} = -5R/12$ and $h_R(0) = 0$ ($\beta_R^{\text{Dirac}} = 0$).
    This changes the relative weight of Dirac vs.\ scalar/vector contributions
    and may shift the ratio from the original $-1/3$.
    If the corrected ratio differs from $-1/3$, update Section~\ref{sec:proven}
    and the prediction file \texttt{prediction\_01\_c1c2\_ratio}.
\end{enumerate}
}

\outcome{
\begin{itemize}
  \item Explicit $h_C^{(s)}$, $h_R^{(s)}$ for scalar and vector fields.
  \item Total SM form factors $F_i^{\text{total}}(z)$ ready for use in
    MR-2 (unitarity), MR-3 (causality), and MR-7 (scattering amplitudes).
  \item Cross-check: local limits reproduce known one-loop $\beta$-functions
    for the full SM.
  \item \textbf{Updated $c_1/c_2$ ratio prediction} with corrected sign
    conventions, ready for experimental confrontation.
\end{itemize}
}

\deps{NT-1 (methodology and conventions established). Uses the same CZ form
factor framework; only traces change.}

\textbf{Status (2026-03-11):} Phases~1--2 \textcolor{proven}{COMPLETE.}
\begin{itemize}
  \item \textbf{Phase 1 (scalar):} $h_C^{(0)}$, $h_R^{(0)}(\xi)$ derived.
    $\beta_W^{(0)} = 1/120$, $\beta_R^{(0)} = \half(\xi - 1/6)^2$.
    142/142 quantitative checks PASS (4 errors caught by independent cross-verification).
  \item \textbf{Phase 2 (vector):} $h_C^{(1)}$, $h_R^{(1)}$ derived.
    $\beta_W^{(1)} = 1/10$, $\beta_R^{(1)} = 0$.
    Ghost counting corrected: 2 Faddeev--Popov ghosts (not 1), shifting
    $\beta_W^{(1)}$ from $13/120 \to 1/10$.
    1519 numerical checks ALL PASS (11 errors caught).
  \item \textbf{Phase 3 (combined SM):} \textcolor{proven}{COMPLETE (2026-03-12).}
    $\alpha_C = 13/120$ (total Weyl${}^2$, $\xi$-independent, parameter-free).
    $\alpha_R(\xi) = 2(\xi - 1/6)^2$ (total $R^2$, depends on Higgs non-minimal coupling).
    $c_1/c_2 = -1/3 + 120(\xi-1/6)^2/13$ (general; $-1/3$ at conformal $\xi=1/6$).
    SM counting: $N_s=4$, $N_D=22.5$, $N_v=12$.
    354/354 quantitative checks ALL PASS.
\end{itemize}

\textbf{Status (2026-03-12):} All three phases \textcolor{proven}{COMPLETE.}
Total: 2015/2015 checks PASS (142 + 1519 + 354).

\subsection{NT-2: Entire-function property of spectral form factors}
\label{task:NT2}

\task{Determine whether the form factors $F_k(\Box/\Lam^2)$ derived from the
spectral action are entire functions of $\Box$, which would explain ghost-free
behavior without additional postulates.}

\method{
\begin{enumerate}[(a)]
  \item From the explicit results of Task~NT-1, analyze the singularity structure
    of $F_k(z)$ in the complex plane.
  \item Verify the Hadamard factorization: if $F_k$ is entire of finite order $\rho$,
    then
    \begin{equation}\label{eq:hadamard}
      F_k(z) = z^m\, e^{P(z)}\, \prod_{n=1}^\infty
      E_p\!\left(\frac{z}{z_n}\right)
    \end{equation}
    with $P$ a polynomial of degree $\leq \rho$ and $E_p$ canonical factors.
    For ghost-freedom, there must be no real positive zeros (no poles in the propagator).
  \item Compare with postulated entire functions in ghost-free gravity
    (Tomboulis, Modesto, Biswas--Mazumdar--Siegel):
    \begin{equation}\label{eq:ghost-free-prop}
      G(k^2) \sim \frac{e^{-\ell(k^2/\Lam^2)}}{k^2 - m^2}
    \end{equation}
    where $\ell$ is an entire function. Determine whether the spectral action
    \emph{derives} this structure.
\end{enumerate}
}

\outcome{Either:
\begin{itemize}
  \item[(A)] $F_k$ are entire $\Rightarrow$ SCT is ghost-free \emph{from first principles},
    a major structural advantage over competing approaches.
  \item[(B)] $F_k$ have isolated poles $\Rightarrow$ identify their physical meaning
    (massive ghost modes, Lee--Wick partners, or artifacts of truncation).
\end{itemize}
}

\deps{Requires NT-1 (explicit form factors).}

\textbf{Status (2026-03-12):} \textcolor{proven}{COMPLETE.}
$\varphi(z)$ entire, Taylor coefficients $a_n = (-1)^n n!/(2n+1)!$.
$F_1(z)$, $F_2(z,\xi)$ entire $\Rightarrow$ ghost-freedom guarantee (outcome~A confirmed).
$\Pi_{\text{entire}}(z) > 0$ for all real $z \geq 0$ (verified to $z=100$).
$F_1(0) = 13/(1920\pi^2)$, $F_2(0, \xi=0) \approx 1.7576 \times 10^{-4}$.
63/63 quantitative checks ALL PASS.
See \texttt{theory/derivations/NT2\_entire\_function.tex}.

\subsection{NT-3: Spectral dimension flow $d_S(\text{scale})$}
\label{task:NT3}

\task{Derive the running spectral dimension $d_S(\ell)$ as a function of
the probing scale $\ell$ from the spectral properties of the Dirac operator $D$.}

\method{
The spectral dimension is defined via the return probability of a diffusion process:
\begin{equation}\label{eq:spectral-dim}
  d_S(\sigma) = -2\,\frac{d \ln P(\sigma)}{d \ln \sigma}, \quad
  P(\sigma) = \Tr\!\left(e^{-\sigma D^2}\right)
\end{equation}
where $\sigma$ is the diffusion time (inverse scale squared).
\begin{enumerate}[(a)]
  \item Compute $P(\sigma)$ from the full spectrum of $D^2$ on the spectral triple.
  \item Expand for large $\sigma$ (IR, should give $d_S \to 4$) and small $\sigma$
    (UV, test for dimensional reduction $d_S \to 2$).
  \item The $4 \to 2$ flow is a universal prediction of many quantum gravity approaches
    (CDT, Ho\v{r}ava--Lifshitz, asymptotic safety). Determine whether SCT reproduces
    this and with what specific flow profile.
\end{enumerate}
}

\outcome{An explicit function $d_S(\sigma)$ interpolating between:
\begin{itemize}
  \item $d_S \to 4$ at large scales (macroscopic spacetime),
  \item $d_S \to d_{\text{UV}}$ at Planckian scales, with $d_{\text{UV}}$
    determined by the spectral triple.
\end{itemize}
If $d_{\text{UV}} = 2$, this connects SCT to the dimensional reduction universality
class. If $d_{\text{UV}} \neq 2$, this is a distinguishing prediction.
}

\deps{Requires the spectral triple to be fully specified (Postulate 5 variants,
see MT-3). A partial result is possible using the minimal spectral triple.}

\textbf{Status (2026-03-16):} \textcolor{proven}{COMPLETE (DEFINITION-DEPENDENT).}
Spectral dimension computed under four definitions:
Propagator (CMN): $d_S = 2$; Heat kernel: $d_S = 4$;
ASZ/fakeon: $d_S = 0 \to 4$; Mittag--Leffler (primary): $d_S \approx 2$ at
ghost scale $\to 4$ in IR.
Critical finding: $P_{\mathrm{ML}}(\sigma) < 0$ for
$\sigma < \sigma_* \approx 0.010\,\Lambda^{-2}$ due to
$\sum R_n = -1.034$ (ghost dominance).
Physical flow: $d_S \approx 2$ at $\sigma \sim 0.05$--$0.1\,\Lambda^{-2}$,
$d_S = 4$ for $\sigma \gg \Lambda^{-2}$.
32/32 verification checks, 2 independent derivations, 5 methods.
See \texttt{theory/derivations/NT3\_spectral\_dimension.tex}.

\subsection{NT-4: Nonlocal gravitational field equations}
\label{task:NT4}

\textbf{Rationale for splitting.} NT-4 is the single largest bottleneck in the
roadmap: 17~downstream tasks depend on it directly or indirectly. The full
nonlinear variational problem on arbitrary backgrounds is substantially harder
than the linearized version on flat spacetime, yet many downstream tasks (MR-2,
MR-3, MR-7, PPN-1) only require the linearized equations. Splitting NT-4 into
three subtasks de-risks the critical path and enables earlier progress on
physical predictions.

\subsubsection{NT-4a: Linearized nonlocal equations on flat background}
\label{task:NT4a}

\task{Linearize the spectral action around flat spacetime
$g_{\mu\nu} = \eta_{\mu\nu} + h_{\mu\nu}$ and derive the linearized field
equations in transverse-traceless gauge.}

\method{
\begin{enumerate}[(a)]
  \item Expand the spectral action~\eqref{eq:nonlocal-eom} to second order
    in $h_{\mu\nu}$. On flat background, $\Box \to \partial^2$ and the form
    factors become ordinary functions of the d'Alembertian:
    $F_i(\partial^2/\Lambda^2)$.
  \item Extract the linearized equations in transverse-traceless gauge.
    The result is a modified wave equation:
    \begin{equation}
      \mathcal{G}_{\mu\nu}^{(1)}[h] +
      \bigl[\alpha_C\,F_1(\partial^2/\Lambda^2) + \alpha_R\,F_2(\partial^2/\Lambda^2)\bigr]
      \,\mathcal{G}_{\mu\nu}^{(1)}[h] = 8\pi G\,T_{\mu\nu}^{(1)}
    \end{equation}
    where the coefficients $\alpha_{C,R}$ arise from the linearized
    Weyl and Ricci contributions.
  \item Verify: (i)~the equations reduce to linearized Einstein equations when
    $F_1, F_2 \to c_1, c_2$; (ii)~gauge invariance under
    $h_{\mu\nu} \to h_{\mu\nu} + \nabla_{(\mu}\xi_{\nu)}$;
    (iii)~the propagator derived from these equations matches MR-2.
  \item Read off the modified Newtonian potential for PPN-1 analysis.
\end{enumerate}
}

\outcome{
\begin{itemize}
  \item Linearized nonlocal field equations ready for propagator analysis (MR-2),
    causality analysis (MR-3), and PPN computation (PPN-1).
  \item Modified Newtonian potential as a function of $\Lambda$ and $F_i$.
  \item Scattering amplitude input for MR-7.
\end{itemize}
}

\deps{NT-1 (Dirac form factors), NT-1b (total form factors for the
physical field equations with correct SM spectrum).}

\textbf{Status (2026-03-12):} \textcolor{proven}{COMPLETE.}
$\Pi_{\text{TT}}(z) = 1 + (13/60)\cdot z\cdot\hat{F}_1(z)$,
$\Pi_s(z,\xi) = 1 + 6(\xi-1/6)^2\cdot z\cdot\hat{F}_2(z,\xi)$.
Scalar decouples at $\xi=1/6$. Barnes--Rivers: $P^{(2)}$ (tr=5), $P^{(0\text{-}s)}$ (tr=1).
Gauge invariance + Bianchi identity verified for 10 random momenta.
Effective masses: $m_2 = \Lambda\sqrt{60/13} \approx 2.148\,\Lambda$,
$m_0(\xi=0) = \Lambda\sqrt{6} \approx 2.449\,\Lambda$.
Modified Newtonian potential:
$V(r)/V_N(r) = 1 - (4/3)\,e^{-m_2 r} + (1/3)\,e^{-m_0 r}$.
88/88 quantitative checks ALL PASS.
See \texttt{theory/derivations/NT4a\_linearized.tex}.

\subsubsection{NT-4b: Full nonlinear variational equations}
\label{task:NT4b}

\task{Derive the full nonlinear nonlocal gravitational field equations from
variation of the spectral action on arbitrary backgrounds:
\begin{equation}\label{eq:nonlocal-eom}
  G_{\mu\nu} + 2\,\frac{\delta}{\delta g^{\mu\nu}}
  \int d^4x\,\sqrt{g}\;\bigl[
  F_1(\Box/\Lambda^2)\,C_{\alpha\beta\gamma\delta}C^{\alpha\beta\gamma\delta}
  + F_2(\Box/\Lambda^2)\,R^2 \bigr] = 8\pi G\, T_{\mu\nu}
\end{equation}
The nontrivial part is the variational calculus of the nonlocal operators
$F_i(\Box)\,\mathcal{R}^2$ with respect to the metric on curved backgrounds.}

\method{
\begin{enumerate}[(a)]
  \item Use the Barvinsky--Vilkovisky technique for variation of nonlocal
    functionals. The key identity is:
    \begin{equation}
      \frac{\delta}{\delta g^{\mu\nu}}\bigl[A\,F(\Box)\,B\bigr]
      = \int_0^1 d\alpha\;\bigl[
      F'(\Box_\alpha)\,(\delta\Box)\,\bigr]_{\alpha\text{-insertion}}
      + \text{boundary terms}
    \end{equation}
    where the $\alpha$-insertion accounts for the non-commutativity of
    $\delta\Box$ with functions of $\Box$.
  \item Specialize to the Weyl tensor and Ricci scalar terms separately.
    For the $R^2$ sector, the variation gives
    $\nabla_\mu\nabla_\nu[F_2(\Box)R]$ terms; for $C^2$, the variation
    involves the Bach tensor and its nonlocal generalization.
  \item Verify: (i)~the equations reduce to Einstein's equations when
    $F_1, F_2 \to c_1, c_2$ (local limit); (ii)~the trace equation is
    consistent; (iii)~the Bianchi identity
    $\nabla^\mu(\text{l.h.s.}) = 0$ is satisfied;
    (iv)~the linearized limit on flat background reproduces NT-4a.
  \item Derive the explicit Wald entropy variation
    $\delta S_{\text{spec}}/\delta R_{\mu\nu\rho\sigma}$ for MT-1.
\end{enumerate}
}

\outcome{
\begin{itemize}
  \item Closed-form nonlinear nonlocal field equations on arbitrary backgrounds.
  \item Wald entropy variation for MT-1.
  \item Foundation for MR-9 (black hole singularity resolution on
    Schwarzschild background).
\end{itemize}
}

\deps{NT-4a (linearized limit as cross-check), NT-1, NT-1b.}

\textbf{Status (2026-03-15):} \textcolor{proven}{COMPLETE.}
Full nonlinear variational field equations derived in the $\{C^2, R^2\}$
(Weyl) basis via the Barvinsky--Vilkovisky alpha-insertion technique.
Key results: field equations with nonlocal Bach tensor $B_{\mu\nu}$,
$H_{\mu\nu}$ from $R^2$ variation, and alpha-insertion corrections
$\Theta^{(C)}$, $\Theta^{(R)}$. Local limit recovers Stelle equations
with $\alpha_C = 13/120$, $\alpha_R(\xi) = 2(\xi-1/6)^2$. Linearized
limit recovers NT-4a propagators. De~Sitter is exact solution. Dual
derivation (Methods~A and~B1) in exact agreement. Wald entropy variation
computed for MT-1. 110~tests, 6/6~consistency checks PASS. See
\texttt{theory/derivations/NT4b\_nonlinear.tex}.

\subsubsection{NT-4c: FLRW reduction of nonlocal equations}
\label{task:NT4c}

\task{Reduce the full nonlocal field equations (NT-4b) to the FLRW metric
$ds^2 = -dt^2 + a(t)^2\,d\vec{x}^2$ and derive the modified Friedmann
equations including nonlocal form factor contributions.}

\method{
\begin{enumerate}[(a)]
  \item Substitute the FLRW metric into the full nonlocal equations of NT-4b.
    On FLRW backgrounds, $\Box = -\partial_t^2 - 3H\partial_t$ (for scalar
    quantities), which simplifies the form factors to functions of the
    conformal time derivative.
  \item Extract the modified Friedmann equations:
    $H^2 + \delta H^2_{\text{NL}}(F_1, F_2) = \frac{8\pi G}{3}\rho$, where
    $\delta H^2_{\text{NL}}$ encodes the nonlocal corrections.
  \item Compute the inflationary perturbation equations to first order in
    $\delta g_{\mu\nu}$ on FLRW background (needed for $n_s$, $r$ in INF-1).
  \item Verify: (i)~GR limit ($F_i \to c_i$) gives standard Friedmann;
    (ii)~de Sitter solution exists; (iii)~consistency with energy conservation
    $\dot{\rho} + 3H(\rho + p) = 0$ (possibly modified).
\end{enumerate}
}

\outcome{
\begin{itemize}
  \item Modified Friedmann equations for cosmological applications (MT-2, INF-1, LT-3c).
  \item Perturbation equations on FLRW for inflationary observables.
  \item De Sitter stability analysis.
\end{itemize}
}

\deps{NT-4b (full nonlinear equations), NT-1b (total form factors).}

\textbf{Status (2026-03-15):} \textcolor{proven}{COMPLETE.}
FLRW reduction of NT-4b field equations derived via independent cross-verification.
Key results: modified Friedmann equations with nonlocal $R^2$ corrections
(Weyl sector drops out: $C_{\mu\nu\rho\sigma} = 0$ on FLRW), GW speed
$c_T = c$ to $\mathcal{O}(H^2/\Lambda^2)$ (GW170817 trivially satisfied),
de~Sitter is exact solution (stable for all~$\xi$), conformal coupling
$\xi = 1/6$ eliminates all corrections, alpha-insertion EOS $w = +1$ (stiff).
Late-time corrections $\mathcal{O}(10^{-120})$.
117/117 pytest tests PASS, multi-method verification, dual derivation agreement.
See \texttt{theory/derivations/NT4c\_flrw.tex}.

\subsection{META-1: Parameter counting and falsifiability criteria}
\label{task:META1}

\task{Perform a systematic parameter count of SCT and establish explicit
falsifiability criteria for each prediction channel.}

\method{
\begin{enumerate}[(a)]
  \item \textbf{Free parameter census.} Enumerate all free parameters of SCT:
    the spectral function $f$ (or equivalently its moments $f_0, f_2, f_4$),
    the finite Dirac operator $D_F$ (13 Yukawa parameters + Majorana masses),
    and the noncommutativity scale $\Lambda$. Determine which are independently
    measurable and which are fixed by internal relations (e.g., sum rules).
  \item \textbf{Comparison with competitors.} Tabulate the parameter count
    against: $\Lambda$CDM (6), SM (19), string theory landscape ($\sim 10^{500}$
    vacua), LQG (Immirzi $\gamma$), asymptotic safety (relevant couplings at
    fixed point).
  \item \textbf{Falsifiability matrix.} For each prediction
    (QNM shifts, $w(z)$, $\Delta F/F$, PPN parameters, $c_1/c_2$,
    inflationary observables, spectral dimension), specify:
    (i)~the predicted value or range as a function of $\Lambda$,
    (ii)~the required experimental precision,
    (iii)~the dataset that tests it,
    (iv)~what outcome would \emph{falsify} SCT (not merely constrain $\Lambda$).
  \item \textbf{Theory review checklist.} Address all 10 mandatory checks
    from the physics methodology: degrees of freedom, action, symmetries,
    observables, domain of validity, limiting cases, predictions, falsification,
    novelty, rigor vs.\ heuristics.
  \item \textbf{Cutoff function $f$ analysis (structural gap).}
    The spectral action $\Tr(f(D^2/\Lambda^2))$ depends on the cutoff function
    $f$, which is an \emph{infinite-dimensional} free input not fixed by any
    postulate. This is the most serious structural ambiguity in the spectral
    action program. Systematically analyze:
    \begin{enumerate}[(i)]
      \item \textbf{Sensitivity audit:} Which predictions depend only on the
        moments $f_0, f_2, f_4$ (i.e., are universal), and which depend on the
        full shape of $f$? The heat-kernel expansion to $a_4$ (Einstein--Hilbert +
        SM Lagrangian) depends only on moments; the nonlocal form factors $F_1$,
        $F_2$ depend on the \emph{entire} function $f$ via the Laplace transform.
      \item \textbf{Classification:} Partition all roadmap predictions into
        ``moment-dependent'' (robust) vs.\ ``$f$-shape-dependent'' (fragile).
        Specifically: $c_1/c_2 = -1/3 + 120(\xi-1/6)^2/13$ is moment-independent
        (universal ratio from heat kernel; reduces to $-1/3$ at conformal coupling
        $\xi = 1/6$); the form factors $F_i(z)$ and their UV behavior
        \emph{are} $f$-dependent.
      \item \textbf{Constraints on $f$:} Determine whether physical requirements
        (unitarity via MR-2, causality via MR-3, finiteness via MR-4/MR-5,
        entire-function property via NT-2) constrain $f$ sufficiently to reduce
        the ambiguity from infinite-dimensional to finite-dimensional.
      \item \textbf{Universality test:} For predictions that do depend on $f$,
        determine whether the \emph{qualitative} predictions (sign, order of
        magnitude, asymptotic behavior) are stable under variation of $f$
        within the allowed class, or whether different choices of $f$ can
        produce qualitatively different physics.
    \end{enumerate}
    This analysis is essential for honestly assessing which SCT predictions
    are genuine and which are artifacts of a particular $f$-choice.
\end{enumerate}
}

\outcome{
\begin{itemize}
  \item Complete parameter table with comparison to competing approaches.
  \item Falsifiability matrix ready for publication.
  \item Self-assessment document identifying what is rigorous, what is
    heuristic, and what remains open in SCT.
  \item Cutoff function sensitivity report classifying all predictions as
    universal (moment-only) vs.\ $f$-dependent, with explicit error bars
    from $f$-variation where applicable.
\end{itemize}
}

\deps{None (meta-task; can proceed immediately). Should be updated
as each prediction task completes.}

% ==========================================================================
\section{\textcolor{mid}{Mid-Term Tasks}}
\label{sec:mid}
% ==========================================================================

These tasks extend SCT to key physical domains. Each requires results from
at least one near-term task.

\subsection{MT-1: Black hole entropy from spectral triple}
\label{task:MT1}

\task{Derive the Bekenstein--Hawking entropy $S_{\text{BH}} = A/(4G)$ from
the spectral triple, and compute leading quantum corrections.}

\method{
\begin{enumerate}[(a)]
  \item \textbf{Entanglement entropy approach.} For a spectral triple
    $(\mathcal{A}, \mathcal{H}, D)$ restricted to a causal diamond bounded
    by a horizon, compute the entanglement entropy of the Dirac vacuum:
    \begin{equation}\label{eq:entanglement}
      S_{\text{ent}} = -\Tr(\rho_{\text{in}}\ln\rho_{\text{in}}), \quad
      \rho_{\text{in}} = \Tr_{\text{out}}|\Omega\rangle\langle\Omega|
    \end{equation}
  \item \textbf{Spectral action on horizon.} Evaluate the spectral action
    on the near-horizon geometry (Rindler approximation) and identify the
    entropy functional:
    \begin{equation}\label{eq:wald-spectral}
      S = -2\pi \int_{\text{horizon}} \frac{\delta S_{\text{spec}}}{\delta R_{\mu\nu\rho\sigma}}
      \, \epsilon_{\mu\nu}\epsilon_{\rho\sigma}\, d\Sigma
    \end{equation}
    (Wald entropy formula applied to the spectral action).
  \item \textbf{Logarithmic corrections.} The spectral action naturally includes
    higher-curvature terms. Compute the correction:
    \begin{equation}\label{eq:log-correction}
      S_{\text{BH}} = \frac{A}{4G} + c_{\log}\,\ln\!\left(\frac{A}{\ell_P^2}\right) + \cdots
    \end{equation}
    and determine $c_{\log}$ from the spectral data. Compare with loop quantum gravity
    ($c_{\log} = -3/2$ for BTZ) and string theory results.
\end{enumerate}
}

\outcome{
\begin{itemize}
  \item Derivation of area law from spectral principles.
  \item Predicted value of $c_{\log}$ as a testable output.
  \item Understanding of how the spectral triple encodes horizon degrees of freedom.
  \item \textbf{All four laws of BH thermodynamics:} zeroth law (constancy of
    surface gravity on horizon), first law (energy balance $dM = (\kappa/8\pi G)\,dA + \ldots$
    with spectral corrections), second law (area non-decrease or generalized
    second law with quantum corrections), third law (extremal limit).
    Note: a 2025 result (arXiv:2509.23423) shows the conventional first law
    is violated in NCG because entropy deviates from Bekenstein--Hawking.
  \item \textbf{Hawking radiation spectrum.} Derive the grey-body factors
    from the nonlocal effective potential around the black hole. The nonlocal
    form factors modify the barrier in the Regge--Wheeler/Zerilli equations,
    potentially altering the emission spectrum at frequencies $\omega \sim \Lambda$.
\end{itemize}
}

\deps{NT-1 (explicit form factors for Wald entropy computation),
NT-1b (total SM form factors for Wald entropy with correct field content),
NT-4b (full nonlinear field equations---variational Wald entropy derivation
requires the complete nonlinear equations, not just the linearized sector).}

\textbf{Note:} NT-3 (spectral dimension) was removed as a dependency.
The near-horizon spectral dimension is an \emph{output} of MT-1 (computed from
the Wald entropy result), not an input. The spectral dimension enters the
\emph{interpretation} of the entropy result, not its derivation.

\textcolor{near}{CONDITIONAL (2026-03-16).}
Wald entropy on Schwarzschild yields
$S = A/(4G) + 13/(120\pi) + (37/24)\ln(A/\ell_P^2) + \mathcal{O}(1)$.
The logarithmic correction $c_{\log} = 37/24$ is a parameter-free prediction
(SM content $+$ graviton via Sen~2012, eq.\,(1.2)).
Four laws: zeroth PASS, first PASS (Iyer--Wald), third PASS,
second \textbf{CONDITIONAL} on MR-2 ghost resolution (Wall stability).
Modified Hawking radiation negligible for astrophysical BH
($T_H/\Lambda \sim 10^{-10}$).
82~tests, 146/146~verification checks PASS across 7~layers.
See \texttt{theory/derivations/MT1\_bh\_entropy.tex}.

\subsection{MT-2: Modified Friedmann equations from SCT}
\label{task:MT2}

\task{Derive the cosmological evolution equations from the spectral action
applied to the Friedmann--Lema\^itre--Robertson--Walker (FLRW) metric.}

\method{
\begin{enumerate}[(a)]
  \item Evaluate the spectral action~\eqref{eq:spectral-action} on the FLRW metric
    $ds^2 = -dt^2 + a(t)^2\,d\vec{x}^2$ with the standard Dirac operator.
  \item Variation with respect to $a(t)$ yields modified Friedmann equations:
    \begin{equation}\label{eq:friedmann-mod}
      H^2 = \frac{8\pi G}{3}\,\rho
      + \delta_1(H/\Lam) + \delta_2(H/\Lam)
    \end{equation}
    where $\delta_1$, $\delta_2$ are corrections from the nonlocal form factors.
  \item Compute the effective equation of state $w_{\text{eff}}(z)$ and compare
    with $\Lambda$CDM and observational data (Planck 2018, DESI DR1/DR2 BAO,
    Pantheon+ SNe, ACT DR6).
  \item Analyze the early-universe regime $H \sim \Lam$: determine whether
    SCT predicts a bounce, inflationary phase, or other modification.
\end{enumerate}
}

\outcome{
\begin{itemize}
  \item Explicit modified Friedmann equations with SCT corrections.
  \item Predicted deviation $w(z) - w_{\Lambda\text{CDM}}(z)$ as function of redshift.
  \item Constraints on $\Lam$ from CMB and BAO data.
  \item Determination of early-universe behavior (bounce vs.\ inflation vs.\ other).
  \item \textbf{Hubble tension benchmark:} Compute $H_0$ from SCT-modified
    expansion history and compare with Planck (indirect, $67.4 \pm 0.5$) and
    SH0ES (direct, $73.0 \pm 1.0$). Maggiore's nonlocal RR model reduces
    $H_0$ tension from $>5\sigma$ to ${\sim}2.0\sigma$ by predicting
    phantom-like $w_{\text{DE}} \approx -1.14$; SCT should benchmark against this.
  \item \textbf{$S_8$ tension benchmark:} Compute $\sigma_8$ from SCT-modified
    growth function $f\sigma_8(z)$ via the modified Poisson equation.
    Nonlocal gravity generically lowers the growth rate, helping alleviate
    the $S_8$ tension (weak lensing vs.\ CMB: KiDS, DES, HSC constraints).
  \item \textbf{Cosmological constant consistency check:} Determine how
    nonlocal form factors modify the effective cosmological constant
    $\Lambda_{\text{CC}}^{\text{eff}}$ relative to the bare spectral action
    value $f_0 \Lambda^4/(16\pi^2)$. Specifically, check whether the nonlocal
    dressing suppresses the bare $\Lambda^4$ contribution (as suggested by
    the spectral action's insensitivity to additive vacuum energy shifts).
    Compare the resulting $\Lambda_{\text{CC}}^{\text{eff}}$ with the causal
    set prediction $\Lambda \sim 1/\sqrt{N}$ (Sorkin's ``everpresent
    $\Lambda$''). This is an internal consistency test: two components of SCT
    (spectral action and causal set sector) predict the same observable via
    different mechanisms, and they must agree or one must dominate.
\end{itemize}
}

\deps{NT-1 (Dirac form factors), NT-1b (total SM form factors---the FLRW
corrections $\delta_1$, $\delta_2$ in Eq.~\eqref{eq:friedmann-mod} depend on
$F_i^{\text{total}}$, not Dirac alone),
NT-2 (entire-function property determines whether cosmological
perturbation theory is well-defined),
NT-4c (FLRW reduction of the nonlocal field equations---the modified
Friedmann equations are obtained by specializing to FRW symmetry).}

\textbf{Status (2026-03-16):} \textcolor{proven}{COMPLETE (NEGATIVE RESULT).}
SCT nonlocal corrections to the Friedmann equation are
$\delta H^2/H^2 = \beta_R(\xi)\,(H/\Lam)^2 \sim 10^{-64}$ at the PPN-1
lower bound ($\Lam = 2.38 \times 10^{-3}$\,eV, $\xi = 0$, $z = 0$), which
is 64~orders of magnitude below the ${\sim}18\%$ needed for the $H_0$~tension.
$w_{\text{eff}} = -1$ to 63~decimal places; $|c_T/c - 1| < 10^{-62}$
(GW170817 by 46~orders); growth modification $\sim 10^{-57}$ ($S_8$ by
55~orders). At $\xi = 1/6$ all corrections vanish. SCT is trivially
consistent with all late-time cosmological data (Planck, DESI, Pantheon+,
KiDS, DES, GW170817). Physical reason: UV nonlocality
($F(\Box/\Lam^2)$, entire functions) cannot generate the IR nonlocality
($\Box^{-1}$, $m \sim H_0$) required for the tensions.
131~tests, 1073~total verification checks (multi-method), 2~figures.
Derivation: \texttt{theory/derivations/MT2\_cosmology.tex}.

\subsection{MT-3: Three variants of Postulate 5}
\label{task:MT3}

\task{Develop, compare, and select among three candidate formulations of
Postulate~5 (the dynamical principle connecting spectral data to spacetime evolution).}

\method{
The three variants under consideration:
\begin{enumerate}[(V1)]
  \item \textbf{Spectral action extremization.}
    Spacetime geometry is determined by extremizing the spectral action:
    \begin{equation}
      \delta\, \Tr\!\left(f\!\left(\frac{D^2}{\Lam^2}\right)\right) = 0
    \end{equation}
    with $D$ depending on the metric (and possibly matter fields) through the
    Dirac operator construction.

  \item \textbf{Spectral flow dynamics.}
    The Dirac operator evolves via a spectral flow equation:
    \begin{equation}
      \frac{\partial D}{\partial t} = \mathcal{F}[D, \text{Spec}(D)]
    \end{equation}
    where $\mathcal{F}$ is a functional determined by the spectral data.
    This is analogous to Ricci flow but for the full spectral triple.

  \item \textbf{Path integral over spectral triples.}
    The quantum theory is defined via a sum over spectral triples:
    \begin{equation}
      Z = \int \mathcal{D}[D]\; e^{-S_{\text{spec}}[D]}
    \end{equation}
    with a suitable measure $\mathcal{D}[D]$ on the space of Dirac operators.
\end{enumerate}

For each variant:
\begin{itemize}
  \item Check mathematical well-posedness (existence, uniqueness, regularity).
  \item Verify that the classical limit gives GR.
  \item Determine what new predictions arise.
  \item Assess compatibility with unitarity and causality.
\end{itemize}
}

\outcome{
\begin{itemize}
  \item Detailed comparison table of three variants.
  \item Identification of the most promising variant (or demonstration that
    they are equivalent in some regime).
  \item Mathematical framework for the chosen variant.
\end{itemize}
}

\deps{NT-1, NT-1b (total SM form factors---assessing well-posedness requires
the physical form factors $F_i^{\text{total}}$, not just the Dirac sector),
NT-2 (needed to assess whether the action principle is well-posed
in each variant).}

\subsection{INF-1: Inflationary predictions from the spectral action}
\label{task:INF1}

\task{Derive the inflationary observables ($n_s$, $r$, $f_{\text{NL}}$,
reheating temperature $T_{\text{reh}}$) from the spectral action's
Starobinsky-like $R^2$ sector, including the effects of the nonlocal
form factors.}

\method{
\begin{enumerate}[(a)]
  \item \textbf{Inflationary potential.} The spectral action's $R^2$ term
    drives Starobinsky inflation. In the Einstein frame, the scalaron
    $\varphi$ has the potential
    \begin{equation}
      V(\varphi) = \frac{3M^2}{4}\left(1 - e^{-\sqrt{2/3}\,\varphi/M_P}\right)^2
    \end{equation}
    with the scalaron mass $M^2 = \Lambda^2/(6 c_2)$ determined by the
    spectral data. Compute $M$ from the total $c_2^{\text{total}}$ (NT-1b).
  \item \textbf{Nonlocal corrections.} The full spectral action replaces
    $c_2 R^2$ with $R\,F_2(\Box/\Lambda^2)\,R$. Following
    Koshelev--Kumar--Starobinsky (JHEP 2020), compute the modified
    slow-roll parameters and their effect on $n_s$, $r$.
    Standard Starobinsky: $n_s \approx 1 - 2/N \approx 0.965$,
    $r \approx 12/N^2 \approx 3.3 \times 10^{-3}$ for $N = 55$ $e$-folds.
    The nonlocal form factors modify $r$ at the few-percent level.
  \item \textbf{Non-Gaussianity.} Compute $f_{\text{NL}}$ from the
    three-point function of curvature perturbations. In single-field
    slow-roll, $f_{\text{NL}} \sim \mathcal{O}(10^{-2})$; the additional
    NCG scalars ($\sigma$-field from seesaw sector, dilaton) could enhance it.
  \item \textbf{Reheating.} In the nonlocal theory, the scalaron's coupling
    to matter fields is modified by the form factors. Compute the decay rate
    $\Gamma_\varphi$ and the resulting reheating temperature $T_{\text{reh}}$.
    Verify consistency with BBN constraints ($T_{\text{reh}} > 5$~MeV).
  \item \textbf{CMB confrontation.} Compare predictions with Planck 2018
    ($n_s = 0.9649 \pm 0.0042$, $r < 0.036$), BICEP/Keck ($r < 0.036$),
    the upcoming Simons Observatory ($\sigma_r \approx 0.003$),
    and LiteBIRD sensitivity ($\sigma_r \approx 10^{-3}$).
  \item \textbf{De Sitter conjecture check.}
    The refined Swampland de Sitter conjecture (Ooguri--Palti--Shiu--Vafa 2018,
    arXiv:1810.05506) states
    $|\nabla V|/V \geq c \sim \mathcal{O}(1)$ or $\min(\nabla_i\nabla_j V) \leq
    -c'\,V$ in any consistent quantum gravity theory. If the spectral action's
    scalar potential (Einstein-frame scalaron + NCG Higgs sector) has a stable
    de Sitter vacuum, determine whether this violates the conjecture and what
    this implies for SCT's relationship to the Swampland program. If the
    potential \emph{does} satisfy the conjecture (quintessence-like behavior),
    determine the constraints on the cutoff function $f$ and $\Lambda$.
    This is a self-consistency check, not a Swampland endorsement---the point
    is to know where SCT stands on a concrete question that competing programs
    (string theory, AS) must also address.
\end{enumerate}
}

\outcome{
\begin{itemize}
  \item Predicted values of $n_s$, $r$, $f_{\text{NL}}$, $T_{\text{reh}}$
    as functions of $\Lambda$ and the spectral data.
  \item Assessment of whether SCT-Starobinsky inflation is within
    LiteBIRD detection range ($r \approx 4 \times 10^{-3}$ at $\sim 4\sigma$).
  \item Comparison with standard Starobinsky, $\alpha$-attractor models,
    and nonlocal Starobinsky (Koshelev et al.).
\end{itemize}
}

\deps{NT-1b (total $c_2^{\text{total}}$ for scalaron mass),
NT-2 (analyticity of form factors for nonlocal slow-roll analysis),
NT-4c (FRW reduction of nonlocal field equations for perturbation theory).}

\textcolor{near}{\textbf{Status (2026-03-15): CONDITIONAL.}}
The full independent cross-verification has been executed.
\textbf{Central negative result:} The SM-only one-loop spectral action at
$\xi = 0$ gives $M = \sqrt{24\pi^2}\,M_P \approx 15.39\,M_P$, which is
$1.2 \times 10^6$ times too heavy for inflation.
Confirmed through 6~independent routes at 100-digit precision.
\textbf{Conditional predictions} (if~$M$ fixed externally to
$3.12 \times 10^{13}\,\text{GeV}$):
$n_s = 0.965$, $r = 3.5 \times 10^{-3}$,
$f_{\text{NL}} = -0.016$, $T_{\text{reh}} = 4.9 \times 10^9\,\text{GeV}$.
All compatible with Planck 2018 and BICEP/Keck 2021.
Nonlocal corrections $< 1\%$ for $\Lambda > 10^{14}\,\text{GeV}$.
Resolution paths: sub-Planckian $\Lambda$ (viable),
sub-Planckian $\Lambda + \text{moderate}\;\xi$ (viable),
higher loops (open).
91/91~pytest tests, 148~verification checks, 3~figures.
Derivation: \texttt{theory/derivations/INF1\_inflation.tex}.

\textbf{Note:} MT-2 is \emph{not} a prerequisite: INF-1 derives the inflationary
dynamics \emph{from} the spectral action directly (Einstein-frame scalaron potential),
not from the modified Friedmann equations. MT-2 and INF-1 are parallel tracks that
independently reduce the spectral action to cosmological settings.

\subsection{INF-2: Dilaton inflation from scale-invariant spectral action}
\label{task:INF2}

\task{Investigate whether a dilaton field arising from the scale-invariant
formulation of the spectral action (Chamseddine--Connes 2006,
arXiv:0512169) can drive successful inflation within the SM spectral triple
alone, without BSM extensions.}

\method{
\begin{enumerate}[(a)]
  \item \textbf{Tree-level dilaton potential.} Derive the Einstein-frame
    potential $V(\varphi)$ from the scale-invariant spectral action
    $S = \operatorname{Tr}(f(D^2/\phi^2))$ after Weyl rescaling
    $g_{\mu\nu} \to e^{-2\varphi} g_{\mu\nu}$.
  \item \textbf{Coleman--Weinberg corrections.} Compute the one-loop
    CW potential from all SM fields, checking whether mass eigenvalues
    depend on $\varphi$ or only on Higgs VEV ratios.
  \item \textbf{Conformal anomaly.} Evaluate the trace anomaly coefficients
    $a$, $c$ for the SM and determine whether they generate a local
    dilaton potential on FLRW backgrounds.
  \item \textbf{$R^2$ kinetic structure.} Analyze whether the $R^2$ term
    in the spectral action provides a kinetic term for $\varphi$
    (it does not---it gives a mass term for $R$, not for $\varphi$).
  \item \textbf{Slow-roll analysis.} For any candidate $V(\varphi)$,
    compute $n_s$, $r$, $f_{\text{NL}}$ and compare with Planck/BICEP.
\end{enumerate}
}

\outcome{
\begin{itemize}
  \item Classification of all potential sources of $V(\varphi)$ in
    the SM spectral triple.
  \item Definitive assessment of whether dilaton inflation is viable
    without BSM content.
  \item Comparison with INF-1 (scalaron inflation) results.
\end{itemize}
}

\deps{INF-1 (scalaron mass result), NT-1b (form factors and SM content),
NT-4b (nonlinear field equations for Weyl rescaling).}

\textcolor{long}{\textbf{Status (2026-03-17): NEGATIVE.}}
Dilaton inflation from the SM spectral triple is excluded by
\textbf{five independent structural obstacles}:
\begin{enumerate}[(1)]
  \item \textbf{Flat tree-level potential:} $V(\varphi) = \Lambda_{\text{cc}}$
    (constant), independent of $\varphi$. The Weyl factors from the
    spectral action density and $\Lambda$-scaling cancel exactly:
    $e^{-4\varphi} \cdot e^{+4\varphi} = 1$.
  \item \textbf{$\varphi$-independent mass eigenvalues:} All SM masses
    $m_i^2 = y_i^2 v^2/2$ depend on Yukawa couplings and the Higgs VEV,
    not on the dilaton. The Coleman--Weinberg potential is therefore
    $V_{\text{CW}}(\varphi) = \text{const}$.
  \item \textbf{No anomaly-driven $V(\varphi)$:} On FLRW, $C^2 = 0$.
    The Euler density $E_4 = 24H^2(\dot{H} + H^2) \neq 0$ locally but
    enters only through non-local functionals (derivative couplings),
    not as a local potential. No anomaly-driven $V(\varphi)$ is generated.
  \item \textbf{CW potential shapes excluded:} Generic CW potentials of the
    form $V = \alpha\,\chi^4(\ln(\chi^2/v^2) - 1/2)$ with SM-scale
    $\alpha$ give $r = 8/N \approx 0.145$, excluded by BICEP/Keck 2021
    ($r < 0.036$).
  \item \textbf{$R^2$ gives scalaron kinetic, not dilaton kinetic:}
    The conformal transformation of $c_2 R^2$ yields a massive scalaron
    $\chi$ with $M^2 = \Lambda^2/(6c_2)$, not a massless dilaton kinetic
    term.
\end{enumerate}
Combined with INF-1 ($M_{\text{scalaron}} = 15.4\,M_P$, too heavy):
\textbf{inflation requires BSM extension or external EFT input}.
91/91~pytest tests (D:~39, DR:~40, suite:~40), $\Sigma$~errors:~0, 0~figures.
Derivation: \texttt{theory/derivations/INF2\_dilaton.tex}.

% ==========================================================================
\section{\textcolor{mid}{Theory Consistency Tasks (MR-1 through MR-9)}}
\label{sec:consistency}
% ==========================================================================

These tasks address the fundamental consistency requirements for a UV-complete
quantum gravity theory. A theory that fails any of the CRITICAL tasks below
cannot claim UV-completeness regardless of other results.

\subsection{MR-1: Lorentzian spectral action}
\label{task:MR1}

\task{Define the spectral action in Lorentzian signature and resolve the
tension between the Lorentzian causal structure (Postulate~2) and the
Euclidean spectral action (Postulate~3).}

\method{
\begin{enumerate}[(a)]
  \item Investigate the \textbf{Krein space formulation} of spectral triples
    (Strohmaier, Paschke--Sitarz, van den Dungen--Rennie): replace the
    Hilbert space $\mathcal{H}$ with a Krein space $(\mathcal{K}, J)$ where
    $J$ is a fundamental symmetry encoding the indefinite metric of
    Lorentzian signature.
  \item Define $D^2$ with Lorentzian signature and determine whether
    $\Tr(f(D^2/\Lambda^2))$ is well-defined (the operator $D^2$ is no
    longer positive-definite).
  \item Prove or disprove that the Wick rotation $z_E \to z_L = -z_E$
    applied to the Euclidean form factors (NT-1) yields a valid Lorentzian
    action at \emph{all} curvature scales, not just in the weak-field limit.
  \item Determine whether the causal set structure (Postulate~2) constrains
    the allowed Lorentzian spectral triples, providing a selection principle.
\end{enumerate}
}

\outcome{Either a well-defined Lorentzian spectral action, or a precise
characterization of the domain of validity of the Euclidean formulation.}

\deps{None (foundational). Can proceed in parallel with all other tasks.}

\textbf{Severity: CRITICAL.}

\textbf{Status (2026-03-14):} \textcolor{proven}{COMPLETE.}
Lorentzian form factors $\varphi(-x)$ derived via Wick rotation $z_E \to z_L = -z_E$.
$\Pi_{\text{TT}}(-k^2/\Lambda^2)$ well-defined for all real $k^2$.
Spectral representation established for the Lorentzian propagator.
189/189 quantitative checks ALL PASS.
See \texttt{theory/derivations/MR1\_derivation.tex}.

\subsection{MR-2: Unitarity and stability}
\label{task:MR2}

\task{Prove that the full graviton propagator in the nonlocal SCT action
has no ghost poles (no negative-residue poles on the physical sheet),
verify the optical theorem, and establish Hamiltonian stability.}

\method{
\begin{enumerate}[(a)]
  \item From the total form factors (NT-1b, Eq.~\eqref{eq:total-ff}),
    construct the full momentum-space graviton propagator in the
    transverse-traceless decomposition:
    \begin{equation}\label{eq:propagator}
      G^{-1}_{\mu\nu\rho\sigma}(k^2) = G^{-1}_{\text{EH}}
      + k^4\bigl[F_1^{\text{tot}}(k^2/\Lambda^2)\,\mathcal{P}^{(2)}
      + F_2^{\text{tot}}(k^2/\Lambda^2)\,\mathcal{P}^{(0)}\bigr]
    \end{equation}
  \item Analyze all poles: positions $k^2 = m_n^2$, residues $Z_n$, and
    their signs. Ghost-freedom requires $Z_n > 0$ for all physical poles.
  \item \textbf{Massive spin-2 mode (Stelle ghost).} In the local truncation
    ($F_1 \to c_1$), the Weyl-squared term produces a massive spin-2 ghost
    with $m^2 = 1/(2|c_1|)$. Determine whether the nonlocal form factors
    $F_1(z)$ cure this ghost (cf.\ Tomboulis, Modesto, Biswas--Mazumdar--Siegel).
  \item \textbf{Residue analysis at the $\Pi_{\text{TT}}$ zero (from NT-4a).}
    The linearized TT propagator denominator $\Pi_{\text{TT}}(z) = 1 + \frac{13}{60}\,z\,\hat{F}_1(z)$
    has a zero at $z_0 \approx 2.41$ (i.e., $k^2 \approx 2.41\,\Lambda^2$),
    producing an effective mass $m_2 = \Lambda\sqrt{60/13} \approx 2.148\,\Lambda$.
    The Yukawa coefficient $-4/3$ matches the Stelle ghost pattern.
    \textbf{Critical subtask:} compute $\Pi_{\text{TT}}'(z_0)$ to determine the
    sign of the propagator residue at this pole.
    If $\Pi_{\text{TT}}'(z_0) > 0$, the pole has negative residue (ghost);
    if $\Pi_{\text{TT}}'(z_0) < 0$, positive residue (healthy massive graviton).
    \emph{Note:} This concerns the \emph{linearized} propagator, not the
    entire-function property of the form factors (NT-2), which guarantees
    ghost-freedom at the nonlinear/analytic level. The two results may be
    reconciled if the full nonlinear theory modifies the pole structure
    relative to the linearized truncation.
  \item \textbf{Scalar mode.} The SCT prediction $3c_1 + c_2 = 0$ implies
    $m_0^2 \to \infty$ (scalar mode decouples). Verify this persists in the
    nonlocal theory: $3F_1(z) + F_2(z) \neq 0$ for $z > 0$?
  \item \textbf{Optical theorem.} Compute the imaginary part of the one-loop
    graviton self-energy and verify $\text{Im}\,\Pi(s) \geq 0$ for $s > 0$.
  \item \textbf{Ostrogradski stability.} Analyze the Hamiltonian structure of
    the nonlocal equations of motion. For infinite-derivative theories, the
    Ostrogradski theorem does not apply directly, but stability of the initial
    value problem must be established.
  \item \textbf{Conformal factor problem.} Verify that the spectral action
    with SCT-fixed coefficients $(c_1, c_2)$ is bounded below in the
    conformal sector, resolving the conformal factor instability of Euclidean
    quantum gravity.
  \item \textbf{Positive energy.} Re-examine the Schoen--Yau / Witten positive
    energy theorem for the nonlocal gravitational action.
\end{enumerate}
}

\outcome{A definitive answer on ghost-freedom, unitarity, and stability of
the SCT gravitational sector.}

\deps{NT-1 (Dirac form factors), NT-1b (total form factors), NT-2
(entire-function property), NT-4a (linearized nonlocal field equations---ghost analysis
requires the explicit linearized equations of motion).}

\textbf{Severity: CRITICAL.}

\textbf{Status (2026-03-17):} \textcolor{proven}{CLOSED ($D^2$-quantization, 87--93\% survival probability).}
\begin{itemize}
  \item \textbf{Ghost catalogue ($|z| \leq 100$):} $\Pi_{\text{TT}}$ has 8 zeros:
    2 real + 3 complex conjugate pairs.
    $z_0 = 2.4148$ (spacelike, $R_0 = -0.4931$),
    $z_L = -1.2807$ (timelike physical ghost, $R_L = -0.5378$),
    Type~C pairs at $|z| \approx 34, 59, 85$ with $|R| < 0.01$ (Lee--Wick, negligible).
  \item \textbf{Scalar sector (SS subtask):} \textcolor{proven}{COMPLETE.}
    Full scalar ghost catalogue at general $\xi \neq 1/6$.
    $\Pi_s(\xi=1/6) \equiv 1$ (conformal decoupling, ghost-free).
    At $\xi = 0$: 8 zeros in $|z| \leq 100$, all complex Lee--Wick pairs;
    dominant pair at $z \approx -2.076 \pm 3.184\,i$ ($|R| \approx 0.469$).
    At $\xi = 1$: real Lorentzian ghost at $z \approx -0.232$ ($R \approx -0.968$),
    violates spectral positivity.
    No positive-real Euclidean zeros for any $\xi$.
    PPN $\gamma \to 1$ for all $\xi$ at solar system distances.
    95/95 verification checks PASS + 55/55 test suite PASS.
    See \texttt{theory/consistency-checks/SS\_scalar\_sector.tex}.
  \item \textbf{Conformal factor:} $\alpha_R(\xi) = 2(\xi-1/6)^2 \geq 0$ for all $\xi$ --- resolved.
  \item \textbf{Ostrogradski:} inapplicable to infinite-derivative theory.
  \item \textbf{Conditions for viability:} (1)~fakeon prescription extends to nonlocal (entire-function) propagators;
    (2)~Donoghue--Menezes unstable ghost mechanism applicable to nonlocal theories;
    (3)~Kubo--Kugo operator-formalism objection resolvable.
  \item \textbf{Ghost width:} $\Gamma/m = C_\Gamma \cdot (\Lambda/M_{\text{Pl,red}})^2$,
    $C_\Gamma = 0.06554$, HLZ partial widths 1:3:6 (scalar:Dirac:vector).
  \item \textbf{Dressed propagator:} $\text{Im}[k^2_{\text{pole}}] < 0$ --- ghost decays (6-step sign chain verified).
  \item \textbf{Optical theorem:} $\text{Im}[T(s)] \geq 0$ for all $s > 0$ with fakeon prescription (one-loop verified).
  \item \textbf{One-loop optical theorem (OT):} \textcolor{proven}{CLOSED ($D^2$-quantization).}
    67/67 OT tests PASS; 10 Devil's Advocate attacks survived.
    $\text{Im}[G_{\text{dressed}}(s)] > 0$ for all $s > 0$ (spectral positivity theorem---rigorous proof, not numerical evidence).
    Central charge $C_m = 283/120$ verified; $N_{\text{eff}} = 143.5$ reconciled.
    See \texttt{theory/consistency-checks/OT\_optical\_theorem.tex}.
  \item \textbf{Commutativity of limits (CL subtask):} \textcolor{proven}{COMPLETE.}
    Rigorous proof that $\lim_{N\to\infty} T_{\text{FK}}(N,s) = T_{\text{FK}}(\infty,s)$
    for all $s > 0$ away from poles. Three independent methods (Weierstrass M-test,
    Lebesgue DCT, Cauchy sequence) all confirm absolute and uniform convergence.
    Weierstrass bound $\sum M_n < 5.01\times 10^{-4}$; complex poles contribute $0.32\%$ of
    total amplitude. 43/43 CL tests PASS; 2868/2868 regression PASS.
    MR-2 Condition~1 is now mathematically established.
    See \texttt{theory/consistency-checks/CL\_commutativity.tex}.
  \item \textbf{Entire part constancy (GZ subtask):} \textcolor{proven}{COMPLETE.}
    Rigorous proof that the entire part $g_A(z)$ in the Mittag-Leffler expansion
    of $1/(z\,\Pi_{\text{TT}}(z))$ is a constant: $g_A(z) = -c_2 = -13/60$
    for all $z \in \mathbb{C}$.  Proof via genus bound (order~1, genus~1
    $\Rightarrow$ $g_A$ at most linear) plus positive real axis limit
    (forces linear coefficient to vanish).  Corollary sum rule:
    $\sum R_n/z_n = c_2 = 13/60$ (verified to 99.95\% with 16~poles).
    Physical implication: the SCT propagator has a pure pole decomposition
    (graviton + ghost tower + constant contact term).
    55/55 GZ tests PASS; 2923/2923 regression PASS.
    See \texttt{theory/consistency-checks/GZ\_entire\_part.tex}.
  \item \textbf{Kubo--Kugo resolution (KK subtask):} \textcolor{near}{PARTIALLY RESOLVED.}
    The Kubo--Kugo objection (arXiv:2308.09006) correctly identifies that the KO quartet
    mechanism does \emph{not} apply to propagator ghosts (spin-2 vs.\ spin-1 quantum number mismatch).
    However, under the fakeon prescription, ghost pair production is identically zero
    ($\text{Im}[G_{\text{FK}}] = 0$ at ghost poles to machine precision at dps=100), and the
    objection does not apply.  Lowest pair-production threshold:
    $E_{\text{th}} = 2\Lambda\sqrt{|z_L|} = 2.264\,\Lambda$ (moot under fakeon).
    Resolution A (KO quartet) REJECTED; Resolution B (fakeon) PRIMARY.
    MR-2 Condition~3 partially resolved: physics clear, rigorous all-orders proof for
    infinite-pole propagators remains open.
    64/64 KK tests PASS; independent re-derivation by 5~methods; dps=100 verification.
    See \texttt{theory/consistency-checks/KK\_kugo\_resolution.tex}.
\end{itemize}
7557+ quantitative checks ALL PASS (MR-2 + A3 + GP + OT + CL + GZ + SS + KK subtasks; 2775 original + 4782 $D^2$-quantization).
See \texttt{theory/derivations/MR2\_derivation.tex}.

\subsection{MR-3: Causality in the nonlocal theory}
\label{task:MR3}

\task{Prove that the nonlocal equations of motion from the spectral action
have a well-posed initial value problem and do not produce acausal propagation.}

\method{
\begin{enumerate}[(a)]
  \item Analyze the \textbf{retarded Green's function} of the linearized
    nonlocal equations of motion. Verify that its support is contained in
    the forward light cone.
  \item Check \textbf{Kramers--Kronig relations} for the graviton self-energy:
    \begin{equation}
      \text{Re}\,\Pi(\omega) = \frac{1}{\pi}\,\mathcal{P}\!\int_{-\infty}^{\infty}
      \frac{\text{Im}\,\Pi(\omega')}{\omega' - \omega}\,d\omega'
    \end{equation}
  \item Compare with known results on causality in infinite-derivative gravity
    (Buoninfante, Modesto, Rachwal): classify the SCT form factors according
    to the causal/acausal taxonomy.
  \item Determine whether the \textbf{initial value problem} for the nonlocal
    field equations is well-posed in the sense of Hadamard (existence,
    uniqueness, continuous dependence on initial data).
  \item \textbf{Macrocausality (cluster decomposition).} Nonlocal theories
    with entire-function form factors violate strict cluster decomposition
    at distances $\sim 1/M$. Prove approximate factorization
    (macrocausality) at distances $\gg 1/M$ for the SCT action.
    Krasnikov (2024) established macrocausality for certain nonlocal $\phi^4$
    models; extend to the gravitational case.
\end{enumerate}
}

\outcome{Proof of causality preservation or identification of acausal
regimes and their physical interpretation. Proof of macrocausality
at distances $\gg 1/\Lambda$.}

\deps{NT-1 (form factors), NT-1b (total SM form factors---causality of
the physical retarded propagator requires $F_i^{\text{total}}$),
NT-2 (analyticity structure), NT-4a (linearized nonlocal field equations).}

\textbf{Severity: CRITICAL.}

\textcolor{near}{CONDITIONAL.}
Completed 2026-03-15 with 83 tests PASS (3061 full regression).
All five subtasks pass:
\begin{itemize}[nosep]
  \item (a)~Retarded G.F.: $G_{\mathrm{ret}}$ vanishes for spacelike separations. PASS.
  \item (b)~Kramers--Kronig: trivially satisfied at tree level ($\mathrm{Im}[\Pi_{\mathrm{TT}}] = 0$). PASS.
  \item (c)~Classification: SCT is Class~III (entire with zeros); closest
    comparison is Anselmi's fakeon gravity.
  \item (d)~IVP: well-posed (classicized, Anselmi--Calcagni 2025, 2 initial data). PASS.
  \item (e)~Macrocausality: $v_{\mathrm{front}} = c$; acausal effects decay
    as $\exp(-1.13\,\Lambda\,r)$. PASS.
  \item Microcausality: VIOLATED at scale $\sim 0.884/\Lambda$ (inherent to fakeon
    prescription; unobservable for $\Lambda \gg$ accessible energies).
\end{itemize}
Condition for upgrade: loop-level Kramers--Kronig verification (requires 2-loop computation).
Independent re-derivation (MR3-DR) confirmed all results with zero discrepancies.
See \texttt{theory/consistency-checks/MR3\_causality.tex}.

\subsection{MR-6: Convergence of the curvature expansion}
\label{task:MR6}

\task{Determine whether the curvature expansion of the spectral action
converges, and if so, with what radius of convergence.}

\method{
\begin{enumerate}[(a)]
  \item Estimate the growth rate of Seeley--DeWitt coefficients $a_{2k}$
    using known results and combinatorial bounds.
  \item Analyze the Laplace transform representation
    $\Tr(f(D^2/\Lambda^2)) = \int_0^\infty \hat{f}(t)\,\Tr(e^{-tD^2/\Lambda^2})\,dt$
    and determine whether it provides a rigorous non-perturbative definition.
  \item For specific spectral triples (e.g., round $S^4$ with radius $a$),
    compute the spectral action \emph{exactly} from the known spectrum of
    $D^2$ and compare with the asymptotic expansion truncated at various orders.
  \item Determine the domain $\|R\|/\Lambda^2 < r_{\text{conv}}$ where the
    $\mathcal{O}(\mathcal{R}^2)$ truncation is reliable.
\end{enumerate}
}

\outcome{Rigorous statement about the relationship between the perturbative
(heat kernel) expansion and the exact spectral action. Radius of convergence
or proof of asymptotic (non-convergent) nature.}

\deps{None (can proceed in parallel with all other tasks).}

\textbf{Severity: IMPORTANT.}

\textbf{Status (2026-03-16):} \textcolor{proven}{COMPLETE.}
Curvature expansion ASYMPTOTIC (Gevrey-1).
Borel radius $\sim 84$.
Laplace representation verified to $168{+}$ digits.
SCT well-defined non-perturbatively via entire form factors.
$K^* = 8$ at $\Lambda^2 a^2 = 1000$.
Reliable for $R/\Lambda^2 < 10^{-10}$ (truncation error $< 10^{-999}$).
74 tests PASS, 4 figures.
See \texttt{MR6\_convergence.tex}.

\subsection{MR-4: Two-loop effective action}
\label{task:MR4}

\task{Compute the two-loop contribution to the gravitational effective action
from the spectral action and verify that no new divergences arise that are
not absorbed by the spectral function $\psi$.}

\method{
\begin{enumerate}[(a)]
  \item Use the background field method with the nonlocal one-loop action
    (NT-1, NT-1b) as the starting point. Expand the metric as
    $g_{\mu\nu} = \bar{g}_{\mu\nu} + h_{\mu\nu}$ and compute the
    two-loop diagrams with nonlocal vertices.
  \item Identify the structure of two-loop divergences: are they of the
    form $R^3$, $\nabla^2 R^2$, or new nonlocal structures?
  \item Determine whether the spectral function $\psi(u)$ (entering via
    $f(D^2/\Lambda^2) = \int_0^\infty \psi(u)\,e^{-uD^2/\Lambda^2}\,du$)
    absorbs all two-loop divergences, or whether new counterterms are needed.
  \item Compare with known two-loop results in higher-derivative gravity
    (Goroff--Sagnotti for pure gravity, van de Ven for gravity + matter).
\end{enumerate}
}

\outcome{Either: (A) all two-loop divergences are absorbed by $\psi$
$\Rightarrow$ evidence for all-orders finiteness, or (B) new counterterms
required $\Rightarrow$ identifies what additional structure SCT needs.}

\deps{NT-1 (one-loop form factors), NT-1b (full SM spectrum), MR-2
(ghost-freedom needed to define loop expansion consistently),
MR-7 (one-loop scattering amplitudes provide the integrand structure
that the two-loop calculation must iterate).}

\textbf{Severity: CRITICAL.}

\textcolor{proven}{UNCONDITIONAL (via CHIRAL-Q).}
The SCT propagator scales as $1/k^2$ in the UV (same as GR, not $1/k^4$ as Stelle).
The $R^3$ Goroff--Sagnotti counterterm \emph{is} generated at two loops, but is
absorbed \emph{unconditionally} by a deformation of the spectral function
$\psi(u) \to \psi(u) + c_2(2-4u+u^2)e^{-u}$
that preserves $f_2$ (Einstein--Hilbert) and $f_4$ (one-loop form factors) while adjusting
$f_6 = \Gamma(3) = 2$.  The CHIRAL-Q chirality theorem (quartic Weyl structure)
establishes unconditional absorption: tensor structure match with $a_6$ is guaranteed
by the chiral block-diagonality of the curvature endomorphism.
Two-loop corrections are parametrically suppressed by $(\Lambda/M_{\text{Pl}})^4/(8\pi^2)^2$
(negligible at all sub-Planckian scales).
Fakeon consistency at two loops follows from Anselmi's theorem (arXiv:2203.02516).
71~tests PASS, 7~independent 120-digit verifications.
See \texttt{MR4\_two\_loop.tex}.

\subsection{MR-5: Power-counting and all-orders finiteness argument}
\label{task:MR5}

\task{Establish a power-counting argument (or prove its impossibility)
for the finiteness of the spectral action at all loop orders.}

\method{
\begin{enumerate}[(a)]
  \item Analyze the superficial degree of divergence $D_\gamma$ for
    $L$-loop diagrams with nonlocal vertices from the spectral action.
    The nonlocal form factors modify the standard power-counting:
    \begin{equation}
      D_\gamma = 4 - E_{\text{ext}} + 2L - \sum_v \delta_v(k^2/\Lambda^2)
    \end{equation}
    where $\delta_v$ encodes the UV suppression from form factors at vertex $v$.
  \item For Schwartz-class spectral functions ($\psi \in \mathcal{S}(\mathbb{R}_+)$),
    determine whether the exponential UV suppression renders all diagrams finite.
  \item For merely exponential spectral functions ($\psi = e^{-u}$), determine
    whether the $1/z$ UV decay of form factors is sufficient for finiteness
    or whether it fails at sufficiently high loop order.
  \item If an all-orders argument exists, formalize it as a theorem with
    precise conditions on $\psi$.
  \item Compare with known results: Tomboulis's argument for entire-function
    gravity, Krasnikov's super-renormalizability claims, and Modesto's
    power-counting for nonlocal gravity.
  \item \textbf{Naturalness and hierarchy problem.} If SCT is UV-finite,
    analyze whether the Higgs mass is protected from quadratic divergences
    (the hierarchy problem). The spectral action's sum rule
    $m_H^2 \propto \Tr(D_F^4) - \Tr(D_F^2)^2/\Tr(\text{Id})$ fixes
    $m_H$ at the unification scale; determine whether the nonlocal form
    factors stabilize this prediction under radiative corrections without
    fine-tuning.
\end{enumerate}
}

\outcome{Either: (A) a theorem stating conditions on $\psi$ for all-orders
finiteness (the strongest UV-completeness result), or (B) identification of
the loop order at which finiteness fails and what is needed to fix it.
Additionally: assessment of naturalness---whether SCT radiative corrections
preserve the Higgs mass prediction without fine-tuning.}

\deps{NT-1 (form factor structure), NT-1b (total SM form factors---power-counting
must account for full matter content and total UV behavior of vertices),
NT-2 (entire-function property),
MR-4 (two-loop as explicit check of any general argument).}

\textbf{Severity: CRITICAL.}

\textbf{Status (2026-03-17):} \textcolor{near}{CONDITIONAL (BV-3,4 only). UV-FINITE in $D^2$-quantization.}
\begin{itemize}
  \item \textbf{All-orders finiteness (Option~A):} NOT achievable by known methods.
    The propagator $\Pi_{\mathrm{TT}} \to -83/6$ gives $G \sim 1/k^2$ (GR-like,
    not Stelle-like $1/k^4$). Superficial degree of divergence $D = 2L+2-E$ grows
    with loop order. Tensor structure mismatch: $N(L)$ invariants vs.\ 1 parameter
    from~$\psi$.
  \item \textbf{Perturbative finiteness (Option~C):} Gevrey-1 asymptotic series
    with optimal truncation at $L_{\mathrm{opt}} \approx 78$ at Planck scale
    ($\varepsilon = 1/(8\pi^2) \approx 0.0127$). Error at optimal truncation
    $\sim e^{-79} \approx 5 \times 10^{-35}$ (Dingle).
  \item \textbf{Improvement over GR:} $39\times$ at Planck scale
    ($L_{\mathrm{opt}} = 78$ vs.\ $L_{\mathrm{break}}^{\mathrm{GR}} = 2$,
    Goroff--Sagnotti).
  \item \textbf{Borel radius:} $R_B^{\mathrm{loop}} \approx 79$, close to
    $R_B^{\mathrm{curv}} \approx 84$ from MR-6 (ratio $\approx 0.94$).
  \item \textbf{Two-loop $D=0$ (MR-5b):} \textcolor{proven}{UNCONDITIONAL (via CHIRAL-Q Theorem~6.12).}
    The perturbative on-shell reduction shows $R_{\mu\nu}\sim\mathcal{O}(\alpha_C)$,
    so at $\mathcal{O}(\alpha_C^2)$ only the CCC invariant
    ($C_{\mu\nu}^{\;\;\kappa\lambda}C_{\kappa\lambda}^{\;\;\rho\sigma}
    C_{\rho\sigma}^{\;\;\mu\nu}$) survives.  With 1~constraint and 1~parameter
    ($\delta f_6$), absorption is unconditional.
    SM CCC coefficient: $-1481/6\approx -246.83$.
    CHIRAL-Q Theorem~6.12 establishes $D=0$ unconditionally via the chirality
    structure of the Seeley--DeWitt coefficients.
    74~pytest PASS (D), 49~self-tests (DR), 34~V-checks at 120~digits.
    See \texttt{theory/derivations/MR5b\_two\_loop.tex}.
  \item \textbf{Key result:} Background field $D = 0$ at $L = 1$ VERIFIED (MR-7).
    Physical $D=0$ at $L=2$ is UNCONDITIONAL (MR-5b via CHIRAL-Q Theorem~6.12):
    on-shell CCC absorption established by chirality structure.
  \item \textbf{Spectral moments:} $f_{2k} = (k-1)!$ for $k = 1,\ldots,10$
    verified at 150-digit precision.
  \item 79/79 pytest tests PASS, 7/7 independent verification checks PASS
    at 150~digits. 5~independent re-derivation methods.
    See \texttt{theory/derivations/MR5\_finiteness.tex}.
  \item \textbf{FUND Program (2026-03-17):} Five fundamental investigations
    executed through a full L$\to$LR$\to$D$\to$DR$\to$V$\to$V3$\to$V4 pipeline,
    probing all known escape routes from the three-loop obstruction:
    \begin{enumerate}[(i)]
      \item \textbf{FUND-FRG} (asymptotic safety): no quantitative connection
        established. $\alpha_C = 13/120$ is $5$--$20\times$ larger than typical
        AS fixed-point values. SM matter content inside extended DEP bound.
      \item \textbf{FUND-FK3} (fakeon at $L=3$): \textcolor{long}{CERTIFIED NEGATIVE.}
        Anselmi's theorem confirms fakeon does not change divergent parts.
        Full Molien series gives $P(4)=3$ independent parity-even quartic
        Weyl invariants (corrected from~2), yielding $3{:}1$ overdetermination.
        All $L\geq 3$ are overdetermined with growing ratio.
      \item \textbf{FUND-NCG} (spectral function constraints):
        \textcolor{long}{CERTIFIED NEGATIVE.}
        NCG axioms do not constrain the spectral function~$f$ (30-year consensus).
        Entropy function (Chamseddine--Connes--van Suijlekom 2018) gives one number
        vs.~three invariants. Non-perturbative correction $\sim e^{-c\Lambda^2}$
        cannot cancel power-law divergences.
      \item \textbf{FUND-SYM} (hidden symmetry): \textcolor{near}{CERTIFIED WITH CORRECTIONS.}
        Cayley--Hamilton identity reduces quartic Weyl invariants from~3 to~2.
        $\Tr(\Omega^4_{\mathrm{chain}})$ gives $(C^2)^2{:}({}^*CC)^2 = 1{:}1$
        (exact, from chiral block-diagonality). However, full~$a_8$ includes
        $[\Tr(\Omega^2)]^2$ with nonzero $pq$ cross-term$\Rightarrow$ ratio
        NOT~1:1 for the full coefficient.
      \item \textbf{FUND-LAT} (lattice/spectral sum): \textcolor{near}{CERTIFIED (informative).}
        Exact spectral action on~$S^4$: $\Delta = 41/15120\times\Lambda^{-4}$
        non-perturbative floor verified to $4\times 10^{-15}$ relative precision.
        $C=0$ on~$S^4$ limits utility for Weyl-sector physics.
    \end{enumerate}
    Revised survival probability: $87$--$93\%$ (post-CHIRAL-Q $D^2$-quantization).
    See \texttt{theory/consistency-checks/FUND\_program.tex}.
\end{itemize}

\subsection{MR-7: Scattering amplitudes and gauge-fixing}
\label{task:MR7}

\task{Compute graviton--graviton scattering amplitudes from the nonlocal
spectral action and establish a consistent gauge-fixing procedure for
the nonlocal theory.}

\method{
\begin{enumerate}[(a)]
  \item \textbf{Gauge-fixing.} The nonlocal action requires a modified
    gauge-fixing procedure. The standard de~Donder gauge
    $\mathcal{F}_\mu = \nabla^\nu h_{\mu\nu} - \half\nabla_\mu h$
    must be supplemented with nonlocal gauge-fixing terms:
    \begin{equation}
      S_{\text{gf}} = \frac{1}{2\alpha}\int d^4x\,\sqrt{g}\;
      \mathcal{F}_\mu\,\mathcal{G}^{\mu\nu}(\Box/\Lambda^2)\,\mathcal{F}_\nu
    \end{equation}
    where $\mathcal{G}^{\mu\nu}$ is chosen to simplify the propagator.
    Determine the optimal choice and verify BRST invariance.
  \item \textbf{Ghost action.} Derive the Faddeev--Popov ghost action for
    the nonlocal gauge-fixing. Verify that ghosts do not introduce new
    physical poles.
  \item \textbf{Tree-level amplitudes.} Compute the $2 \to 2$ graviton
    scattering amplitude at tree level with the full nonlocal propagator.
    Verify that: (i) the amplitude reduces to GR in the IR, (ii) the UV
    behavior is improved, (iii) no acausal features appear.
  \item \textbf{One-loop amplitudes.} Compute the one-loop $2 \to 2$
    amplitude and verify: (i) gauge invariance (Ward identities),
    (ii) unitarity (optical theorem), (iii) UV finiteness.
  \item \textbf{Matter sector unitarity.} Formulate and prove a theorem that the spectral
    invariance of $\Tr(f(D^2/\Lambda^2))$ automatically guarantees the Ward--Takahashi
    and Slavnov--Taylor identities for graviton--fermion and graviton--gauge boson scattering,
    ensuring that the nonlocal form factors do not spoil matter sector unitarity.
\end{enumerate}
}

\outcome{
\begin{itemize}
  \item Consistent gauge-fixed nonlocal action with BRST symmetry.
  \item Explicit graviton scattering amplitudes at tree and one-loop level.
  \item Verification of unitarity, gauge invariance, and improved UV behavior.
\end{itemize}
}

\deps{NT-1b (total form factors), NT-4a (linearized nonlocal field equations---needed
for the external states in scattering amplitudes),
MR-2 (ghost-freedom of propagator),
MR-3 (causality of Green's functions).}

\textbf{Severity: IMPORTANT.}

\textbf{Status (2026-03-15):} \textcolor{proven}{COMPLETE.}
Tree-level amplitude verified identical to GR by Modesto--Calcagni field
redefinition theorem (2107.04558); all three conditions satisfied ($E \cdot F \cdot E$
form with $V=0$, entire form factors, vacuum background). One-loop UV behavior:
logarithmic ($D=0$ for graviton bubble, same as Stelle gravity), with counterterms
$\delta\alpha_C\, C^2 + \delta\alpha_R\, R^2$. Ward identity $k^\mu P^{(2)}_{\mu\nu,\rho\sigma}=0$
verified to machine precision. FP ghost sector in minimal de~Donder gauge: standard
$1/k^2$ propagator, no new poles.
Verified by independent cross-verification (literature review, dual derivation, numerical checks),
64 derivation tests, 5 independent re-derivation methods, 10 verification checks at 120-digit precision.
See \texttt{theory/derivations/MR7\_scattering.tex}.

\subsection{MR-8: Non-perturbative definition}
\label{task:MR8}

\task{Define the quantum theory non-perturbatively via a path integral
(or equivalent construction) over Dirac operators and establish the
existence and properties of the measure $\mathcal{D}[D]$.}

\method{
\begin{enumerate}[(a)]
  \item Formalize the space $\mathcal{M}_D$ of Dirac operators compatible
    with the spectral triple axioms. Determine its topology and geometry.
  \item Define a measure $\mathcal{D}[D]$ on $\mathcal{M}_D$ that is:
    (i) diffeomorphism-invariant, (ii) positive-definite (Euclidean) or
    appropriately oscillatory (Lorentzian), (iii) compatible with the
    spectral action as the weight function.
  \item Investigate discretized approximations: the spectral action on
    finite-dimensional spectral triples (matrix geometries, fuzzy spaces)
    as regularized versions of the full path integral.
  \item Compare with: Causal Dynamical Triangulations (CDT), Euclidean
    Dynamical Triangulations (EDT), and random matrix models for
    noncommutative geometries (Barrett--Glaser).
  \item \textbf{Causal set convergence (background independence check).}
    For causal sets obtained by Poisson sprinkling on known manifolds
    ($S^4$, $T^4$, Schwarzschild), construct the truncated spectral triple
    (finite matrix $D_N$) following the Glaser--Stern method (2020--2021)
    and verify convergence of the spectral action
    $\Tr(f(D_N^2/\Lambda^2)) \to \Tr(f(D^2/\Lambda^2))$ as $N \to \infty$.
    This validates the assumption that SCT's continuum calculations (NT-1,
    NT-1b, NT-4a/b) are the correct large-$N$ limit of the fundamental discrete
    structure postulated in Postulates~1--2.
  \item \textbf{Benincasa--Dowker action limit.}
    The Benincasa--Dowker (BD) action is the \emph{only} known discrete action
    on causal sets that reproduces $\int R\sqrt{g}\,d^4x$ in the continuum limit
    (Benincasa--Dowker 2010, arXiv:1001.2725). Verify that the spectral action
    on Poisson sprinklings, $\Tr(f(D_N^2/\Lambda^2))$, reduces to the BD action
    (or a known generalization) in the appropriate regime. If it does not,
    identify the additional structure in the spectral action and determine
    whether this is a feature (richer discrete dynamics) or a bug (inconsistency
    with established causal set results). This is a key consistency check for
    Postulate~2.
\end{enumerate}
}

\outcome{A rigorous (or at least well-motivated) non-perturbative definition
of the SCT quantum theory, with evidence from finite-dimensional models.
In particular, numerical evidence that truncated spectral triples on Poisson
sprinklings converge to continuum spectral geometry.}

\deps{MT-3 (Postulate~5 variant V3 provides the starting point), MR-1
(Lorentzian vs.\ Euclidean formulation).}

\textbf{Severity: IMPORTANT.}

\subsection{MR-9: Black hole singularity resolution}
\label{task:MR9}

\task{Determine whether the nonlocal SCT equations of motion resolve
the Schwarzschild singularity, producing a regular black hole geometry.}

\method{
\begin{enumerate}[(a)]
  \item Solve the linearized nonlocal field equations
    $G_{\mu\nu} + \Lambda^{-2}[\text{nonlocal terms}] = 0$ for a
    spherically symmetric static ansatz $ds^2 = -A(r)\,dt^2 + B(r)\,dr^2
    + r^2\,d\Omega^2$.
  \item Determine whether the form factors $F_1$, $F_2$ produce a
    de~Sitter core at $r \to 0$ (as in Modesto's nonlocal gravity),
    a bounce, or some other regular structure.
  \item Compute the Kretschmer scalar $R_{\mu\nu\rho\sigma}R^{\mu\nu\rho\sigma}$
    and verify it remains finite everywhere.
  \item Compare with known regular black hole solutions: Hayward, Bardeen,
    Frolov, and the Modesto--Bambi nonlocal solution.
  \item If regular: determine the mass-radius relation for the de~Sitter
    core and any observable signatures (modified Hawking temperature,
    quasi-normal modes).
  \item \textbf{Penrose singularity theorem status.} Determine which
    energy condition is violated by the nonlocal corrections, and whether
    the geodesic focusing theorem (Raychaudhuri equation) acquires
    defocusing terms from the SCT form factors. This explains
    \emph{why} singularity resolution occurs at the level of the
    singularity theorems, not just in specific solutions.
\end{enumerate}
}

\outcome{Either: (A) SCT resolves the singularity, providing a concrete
prediction for the interior geometry and identification of which Penrose
theorem premise fails, or (B) the singularity persists,
indicating the $\mathcal{O}(\mathcal{R}^2)$ truncation is insufficient.}

\deps{NT-1 (Dirac form factors), NT-1b (total SM form factors---the physical
BH field equations require $F_i^{\text{total}}$, not just the Dirac sector),
NT-2 (entire-function property---needed for singularity resolution mechanism),
NT-4b (full nonlinear field equations for spherical symmetry),
MT-1 (black hole entropy for consistency check).}

\textbf{Severity: DESIRABLE.}

\subsection{PPN-1: Parameterized post-Newtonian parameters}
\label{task:PPN1}

\task{Compute the full set of parameterized post-Newtonian (PPN) parameters
for the nonlocal spectral action and confront with Solar System tests.}

\method{
\begin{enumerate}[(a)]
  \item \textbf{Linearized potential.} From the linearized nonlocal field equations (NT-4a),
    solve for the static spherically symmetric weak-field metric
    $ds^2 = -(1 + 2\Phi)\,dt^2 + (1 - 2\Psi)\,d\vec{x}^2$ with the nonlocal
    form factors. For exponential-type form factors, the Newtonian potential is
    \begin{equation}
      \Phi(r) = -\frac{Gm}{r}\,\text{erf}\!\left(\frac{Mr}{2}\right)
    \end{equation}
    where $M$ is the mass scale set by $\Lambda$. This implies
    $\gamma_{\text{PPN}} \to 1$ and $\beta_{\text{PPN}} \to 1$ at distances
    $r \gg 1/M$ with exponentially suppressed corrections.
  \item \textbf{Full PPN formalism.} Extend to the full 10-parameter PPN framework
    (Will 2014). Compute all velocity-dependent and preferred-frame parameters
    ($\alpha_1, \alpha_2, \alpha_3, \zeta_1, \ldots$). The spectral action is
    spectrally invariant (hence diffeomorphism-invariant), which should set all
    preferred-frame parameters to zero.
  \item \textbf{Confrontation with data:}
    \begin{itemize}
      \item Cassini bound: $|\gamma - 1| < 2.3 \times 10^{-5}$
      \item Lunar Laser Ranging: Nordtvedt effect, $|\eta| < 4.4 \times 10^{-4}$
      \item Perihelion precession of Mercury: constrains $\beta$
      \item Binary pulsars (Hulse--Taylor, double pulsar): radiation damping
    \end{itemize}
  \item \textbf{Short-distance deviations.} Compute $\Delta F/F(r)$ at
    sub-millimeter scales for comparison with E\"ot--Wash torsion balance
    experiments and atom interferometry. Determine the lower bound on $\Lambda$
    from null results.
\end{enumerate}
}

\outcome{
\begin{itemize}
  \item Full PPN parameter table for SCT as a function of $\Lambda$.
  \item Lower bound on $\Lambda$ from Solar System and laboratory tests.
  \item Demonstration that SCT passes all current weak-field tests.
  \item This is the single most important calculation for experimental
    credibility: without PPN parameters, SCT cannot be placed in the
    standard framework used by experimentalists to compare gravitational
    theories.
\end{itemize}
}

\deps{NT-4a (linearized nonlocal field equations---the weak-field limit requires
the explicit linearized equations), NT-1b (total form factors---the PPN
parameters depend on $F_i^{\text{total}}$).}

\textbf{Severity: CRITICAL.} This is the gold standard for testing
gravitational theories.

\textbf{Status (2026-03-14):} \textcolor{proven}{COMPLETE.}
\begin{itemize}
  \item $\gamma_{\text{PPN}} = 1 + \tfrac{2}{3}\,e^{-m_2 r} + O(e^{-m_0 r})$ (Level~2 Yukawa, derived).
  \item At $r = 1\,\text{AU}$, $\Lambda \geq 2.38\times 10^{-3}\,\text{eV}$:
    $|\gamma - 1| \sim e^{-10^{14}} \approx 0$ (indistinguishable from GR).
  \item SCT passes ALL solar system tests trivially (Cassini, LLR, perihelion, binary pulsars).
  \item All preferred-frame parameters zero (diffeomorphism invariance of spectral action).
\end{itemize}
2829/2829 quantitative checks ALL PASS.
See \texttt{theory/predictions/prediction\_ppn1.tex}.

\subsection{FUND-1: Recovery of QFT on curved spacetime}
\label{task:FUND1}

\task{Demonstrate that SCT reduces to standard quantum field theory on
curved spacetime in the appropriate limit.}

\method{
\begin{enumerate}[(a)]
  \item \textbf{Semiclassical limit.} Starting from the full SCT path integral
    (or its perturbative expansion), perform a Born--Oppenheimer-type separation
    of gravitational and matter degrees of freedom: fix the gravitational sector
    to a classical background $\bar{g}_{\mu\nu}$ and show that the matter sector
    reduces to QFT on $(\mathcal{M}, \bar{g}_{\mu\nu})$.
  \item \textbf{Algebraic QFT structure.} Verify that in the regime of weak
    and slowly varying gravitational fields, the matter algebra satisfies
    the Haag--Kastler axioms (isotony, locality, covariance) to the extent
    expected of QFT on curved spacetime.
  \item \textbf{Stress-energy renormalization.} Show that the renormalized
    stress-energy tensor $\langle T_{\mu\nu}\rangle_{\text{ren}}$ in the
    semiclassical limit reproduces the standard conformal anomaly:
    $\langle T^\mu{}_\mu\rangle = a\,E_4 + c\,C^2 + b\,\Box R$ with
    the correct central charges for the SM matter content.
  \item \textbf{Comparison:} The spectral action's heat-kernel expansion
    already reproduces the classical SM Lagrangian coupled to gravity.
    The nontrivial step is showing that the \emph{quantum} theory on
    the background reproduces standard QFT including renormalization,
    anomalies, and running couplings.
  \item \textbf{Fermion doubling specification.} The noncommutative geometry
    approach to the Standard Model encounters the fermion doubling problem
    (Barrett 2006; Connes 2006): the real structure $J$ on the finite spectral
    triple naturally produces twice the physical fermion degrees of freedom.
    The standard resolution (Connes--Chamseddine ``first-order condition'' or
    the ``Majorana condition'' on the full Hilbert space) must be:
    (i)~explicitly stated as part of the SCT formulation,
    (ii)~shown to be compatible with the Lorentzian spectral action (MR-1),
    (iii)~verified to not introduce anomalies or spoil the SM quantum numbers.
    If the resolution changes the fermion content (as in some approaches), the
    effect on form factors (NT-1, NT-1b) must be assessed.
\end{enumerate}
}

\outcome{
\begin{itemize}
  \item Proof that SCT reduces to QFT on curved spacetime in the
    appropriate semiclassical limit.
  \item Explicit conditions on the gravitational background (curvature
    $\ll \Lambda^2$) under which this reduction is valid.
  \item Verification of conformal anomaly coefficients.
\end{itemize}
}

\deps{NT-1b (total SM form factors), NT-4b (full nonlinear field equations for
the background), MT-3 (quantum formulation needed for semiclassical limit).}

\textbf{Severity: IMPORTANT.} This is a fundamental consistency requirement:
it connects the UV-complete theory to the regime where all existing particle
physics experiments operate.

\subsection{COMP-1: Cross-program comparison table}
\label{task:COMP1}

\task{Produce a detailed comparison of SCT with competing quantum gravity
programs on all shared prediction channels.}

\method{
\begin{enumerate}[(a)]
  \item \textbf{Programs:} Loop quantum gravity, asymptotic safety,
    causal dynamical triangulations, string theory (where applicable),
    infinite-derivative gravity (Tomboulis--Modesto--BMS).
  \item \textbf{Comparison axes:} spectral dimension $d_S(\ell)$,
    BH entropy $c_{\log}$, singularity resolution mechanism, inflationary
    predictions ($n_s$, $r$), modified dispersion relations, PPN deviations,
    UV behavior (asymptotically free / safe / finite), cosmological
    constant treatment, matter coupling.
  \item \textbf{Distinguishing predictions.} Identify observables where SCT
    gives quantitatively different predictions from each competitor.
    The ratio $c_1/c_2$ is already unique; determine if there are others.
  \item \textbf{Quantitative asymptotic safety (AS) comparison.}
    Asymptotic safety is the closest competitor to SCT in methodology (both use
    the gravitational effective action with higher-curvature operators). Perform
    a direct, quantitative comparison:
    \begin{enumerate}[(i)]
      \item Compare the AS fixed-point values of $g^* = G\,k^2$,
        $\lambda^* = \Lambda_{\text{cc}}\,k^{-2}$ (Reuter et al.) with the
        SCT spectral action's running couplings $G(\Lambda)$, $\Lambda_{\text{cc}}(\Lambda)$.
      \item Compare the AS prediction for the spectral dimension
        $d_S^{\text{AS}}(\ell) \to 2$ (Lauscher--Reuter 2005) with the SCT
        prediction from NT-3. If both give $d_S \to 2$, this is universality;
        if they differ, this is a distinguishing observable.
      \item Compare the AS predictions for the graviton propagator (Codello et al.\
        2008) with the SCT propagator from NT-4a/MR-2. The form factor structure
        $F_i(z)$ from the spectral action is a concrete prediction that AS
        computes differently (polynomial truncation vs.\ nonlocal entire function).
    \end{enumerate}
  \item Update as each prediction task (MT-1, MT-2, INF-1, LT-3, PPN-1,
    NT-3) is completed.
\end{enumerate}
}

\outcome{
\begin{itemize}
  \item Publication-ready comparison table.
  \item Identification of the most discriminating experiments/observations
    between SCT and competing approaches.
  \item Quantitative AS vs.\ SCT comparison on spectral dimension,
    propagator, and running couplings.
\end{itemize}
}

\deps{Accumulates results from all prediction tasks. Can begin with
existing results and be updated iteratively.}

\textbf{Severity: DESIRABLE.} Important for positioning SCT in the field,
but not a theory-internal consistency requirement.

% ==========================================================================
\section{\textcolor{long}{Long-Term Tasks}}
\label{sec:long}
% ==========================================================================

These are the ultimate goals of the SCT program. Each requires substantial
results from mid-term and consistency tasks.

\subsection{LT-1: UV-completeness}
\label{task:LT1}

\task{Prove the existence of a fixed-point structure (asymptotic freedom,
asymptotic safety, or finiteness) for the full nonlocal spectral action functional
under the renormalization group.}

\method{
\begin{enumerate}[(a)]
  \item Define the RG flow for the spectral action. The natural cutoff is $\Lam$
    appearing in~\eqref{eq:spectral-action}. Define the effective action at scale $k$:
    \begin{equation}
      \Gamma_k = \Tr\!\left(f_k\!\left(\frac{D^2}{k^2}\right)\right)
    \end{equation}
    where $f_k$ flows with the Wetterich--Morris exact RG equation.
  \item Compute the beta functions for the couplings encoded in $f_k$.
    In the truncation to local operators, these include $G(k)$, $\Lam_{\text{cc}}(k)$,
    $c_1(k)$, $c_2(k)$.
  \item Search for fixed points: $\beta_i(g^*) = 0$ for all couplings.
  \item If the entire-function property (NT-2) holds, investigate whether the
    nonlocal structure is preserved under RG flow (``nonlocal asymptotic safety'').
  \item Prove or disprove finiteness: if the spectral action is UV-finite at
    all loop orders, this would establish SCT as a UV-complete theory without
    requiring a fixed point.
  \item \textbf{UV-completeness criterion clarification.}
    Before attempting the proof, establish what ``UV-complete'' means for SCT
    specifically. The three possible outcomes are:
    \begin{enumerate}[(i)]
      \item \textbf{Asymptotic freedom} ($g_i(k) \to 0$ as $k \to \infty$):
        the theory is weakly coupled in the UV. Requires identification of
        the relevant couplings and their one-loop beta functions.
      \item \textbf{Asymptotic safety} ($g_i(k) \to g_i^*$ as $k \to \infty$,
        $g_i^* \neq 0$ with finitely many relevant directions): the theory
        has an interacting UV fixed point with finite predictive power.
      \item \textbf{UV finiteness}: all scattering amplitudes are finite at
        all orders without renormalization. This is the strongest claim and
        would require proof that the nonlocal form factors $F_i$ suppress
        all UV divergences.
    \end{enumerate}
    The criterion must be stated \emph{before} the calculation, not chosen
    post-hoc to match whatever is found.
\end{enumerate}
}

\outcome{
\begin{itemize}
  \item Proof (or strong evidence) for UV-completeness of SCT.
  \item Classification: asymptotically free / safe / finite / none.
  \item If none: identification of what additional structure is needed.
  \item Explicit criterion for UV-completeness stated in advance of the
    calculation, with clear success/failure conditions.
\end{itemize}
}

\deps{NT-1 (Dirac form factors), NT-1b (total SM form factors---the RG beta
functions $\beta(G)$, $\beta(c_1)$, $\beta(c_2)$ depend on the full
matter content), NT-2 (analyticity structure), MT-3 (correct formulation
of quantum theory), MR-2 (unitarity and stability---UV-completeness requires
the theory to be ghost-free), MR-4 (two-loop calculation---determines
subleading UV structure), MR-5 (all-orders finiteness argument).}

\subsection{LT-2: Information paradox and island formula}
\label{task:LT2}

\task{Derive the island formula for entanglement entropy from SCT principles,
and resolve the black hole information paradox within the spectral framework.}

\method{
\begin{enumerate}[(a)]
  \item Start from the entanglement entropy of the spectral triple (MT-1).
  \item Derive the quantum extremal surface (QES) condition from the spectral action:
    \begin{equation}\label{eq:qes}
      \frac{\delta}{\delta X^\mu}\left[
      \frac{\text{Area}(\gamma_A)}{4G_{\text{eff}}} + S_{\text{bulk}}(\Sigma_A)
      \right] = 0
    \end{equation}
    where $G_{\text{eff}}$ is the effective Newton constant from the spectral action
    and $S_{\text{bulk}}$ is the bulk entanglement entropy.
  \item Show that the Page curve emerges from the competition between
    ``no-island'' and ``island'' saddles in the spectral path integral.
  \item Determine whether the spectral triple naturally implements the
    ``code'' model: $S(A) = \text{Area}(\gamma_A)/(4G_{\text{eff}}) + S_{\text{bulk}}$
    as a \emph{definition} of geometry from entanglement data.
\end{enumerate}
}

\outcome{
\begin{itemize}
  \item Island formula derived from spectral principles (not postulated).
  \item Explicit Page curve for evaporating black hole in SCT.
  \item Understanding of how information is preserved in the spectral framework.
  \item \textbf{Firewall resolution.} SCT's inherent nonlocality provides a
    natural resolution mechanism (Giddings-type nonlocal information transfer):
    the form factors spread quantum effects over scales $\sim 1/\Lambda$ near
    the horizon, potentially smoothing the firewall without violating
    monogamy of entanglement. Make this connection explicit and determine
    whether the AMPS paradox is resolved or merely reframed.
\end{itemize}
}

\deps{MT-1 (black hole entropy), MT-3 (quantum formulation), LT-1 (UV-completeness
needed for late-time behavior).}

\subsection{LT-3: Testable numerical predictions}
\label{task:LT3}

The original single task is split into four focused subtasks, each targeting
a distinct experimental channel with its own methodology and dataset.

\subsubsection{LT-3a: Quasinormal mode shifts}
\label{task:LT3a}

\task{Compute the QNM frequency shifts for Schwarzschild and Kerr black
holes from the nonlocal SCT action and confront with gravitational wave
ringdown observations.}

\method{
\begin{enumerate}[(a)]
  \item From the linearized nonlocal field equations (NT-4b) on a Schwarzschild
    background, derive the modified Regge--Wheeler and Zerilli equations with
    nonlocal potentials depending on $F_1(\Box/\Lambda^2)$, $F_2(\Box/\Lambda^2)$.
  \item Compute the QNM frequencies:
    \begin{equation}
      \omega_{\text{QNM}}^{\text{SCT}} = \omega_{\text{QNM}}^{\text{GR}}
      + \delta\omega\!\left(\frac{c_1}{\Lam^2},\, \frac{c_2}{\Lam^2}\right)
    \end{equation}
    with the ratio $c_1/c_2$ fixed by the spectral triple (NT-1b step~e).
    The shifts $\delta\omega$ depend on a single parameter $\Lam$.
  \item Extend to Kerr using the modified Teukolsky equation.
  \item Compare with LVK O4/O5 ringdown data; project sensitivity
    for LISA and Einstein Telescope (ET).
\end{enumerate}
}

\outcome{QNM frequency table as function of $\Lambda$; bound on $\Lambda$
from current ringdown observations; forecast for next-generation detectors.}

\deps{NT-1b, NT-4b (nonlinear equations on BH background), MT-1 (BH background),
MR-9 (for near-singularity QNMs).}

\subsubsection{LT-3b: Gravitational wave propagation}
\label{task:LT3b}

\task{Compute the GW speed, polarization content, and frequency-dependent
dispersion from the nonlocal SCT action.}

\method{
\begin{enumerate}[(a)]
  \item \textbf{GW speed.} From the linearized equations (NT-4a), derive the
    group and phase velocities of gravitational waves on an FLRW background.
    Ghost-free infinite-derivative gravity predicts $c_{\text{GW}} = c$
    (Buoninfante et al.\ 2019), consistent with the GW170817 bound
    $|c_{\text{GW}}/c - 1| < 10^{-15}$. Verify this for SCT's specific
    form factors.
  \item \textbf{Polarization.} Determine the number of propagating polarization
    modes. Ghost-free IDG gives only 2 tensor modes (same as GR); verify this
    persists in the full nonlocal theory with SCT-fixed coefficients.
  \item \textbf{Dispersion.} The nonlocal form factors introduce
    frequency-dependent corrections to the dispersion relation:
    $\omega^2 = k^2 + \delta\omega^2(k;\, \Lambda)$.
    Compute $\delta\omega^2(k)$ and constrain using multi-messenger observations
    (GW170817/GRB~170817A), Fermi-LAT/GBM time-of-flight bounds on
    Lorentz-violating dispersion, and IceCube/KM3NeT high-energy neutrino
    timing constraints.
  \item \textbf{Secondary GW spectrum.} If INF-1 produces enhanced small-scale
    perturbations, compute the induced secondary gravitational wave background.
    Compare with PTA observations (NANOGrav, EPTA, PPTA).
\end{enumerate}
}

\outcome{GW propagation properties of SCT; bounds from multi-messenger
astronomy; compatibility with existing constraints.}

\deps{NT-4a (linearized propagation equations), NT-4c (FLRW background),
INF-1 (for secondary spectrum).}

\subsubsection{LT-3c: Cosmological predictions}
\label{task:LT3c}

\task{Produce numerical predictions for the cosmological equation of state
$w(z)$ and structure formation observables from the SCT-modified Friedmann
equations.}

\method{
\begin{enumerate}[(a)]
  \item From MT-2 (modified Friedmann equations), compute:
    \begin{equation}
      w(z) = -1 + \delta w(z;\, \Lam)
    \end{equation}
    Fit to the CPL parametrization $w(z) = w_0 + w_a\, z/(1+z)$.
  \item Implement SCT modifications in CLASS/hi\_class + cobaya.
    Compute the CMB temperature and polarization power spectra $C_\ell^{TT}$,
    $C_\ell^{EE}$, $C_\ell^{TE}$ and the matter power spectrum $P(k)$.
  \item Run MCMC fits against: Planck 2018 + ACT DR6 (CMB), DESI DR1/DR2 (BAO),
    Pantheon+ / Union3 (SNe Ia), KiDS-1000/DES-Y3 (weak lensing). Report
    constraints on $\Lambda$ and comparison with $\Lambda$CDM (Bayesian evidence,
    $\Delta\chi^2$, AIC/BIC).
  \item Project sensitivity of Euclid, LSST (Vera Rubin), CMB-S4,
    Simons Observatory ($\sigma_r \approx 0.003$).
\end{enumerate}
}

\outcome{Constraints on $\Lambda$ from cosmological data; model comparison
with $\Lambda$CDM; forecast for next-generation surveys.}

\deps{MT-2, NT-1b, NT-4c (FLRW reduction for cosmological perturbations).}

\subsubsection{LT-3d: Laboratory and Solar System tests}
\label{task:LT3d}

\task{Compute the short-distance modification of gravity and confront with
laboratory and Solar System experiments.}

\method{
\begin{enumerate}[(a)]
  \item The nonlocal form factors modify the Newtonian potential at short
    distances:
    \begin{equation}
      V(r) = -\frac{Gm_1 m_2}{r}\left(1 + \frac{\Delta F}{F}(r;\, \Lam)\right)
    \end{equation}
    Compute $\Delta F/F(r)$ from the exact form factors (NT-1b, NT-4a).
  \item Compare with torsion balance experiments (E\"ot--Wash group:
    $\alpha < 10^{-2}$ at $r \sim 50\;\mu$m), atom interferometry,
    and Casimir force measurements.
  \item Combine with PPN-1 results for Solar System constraints.
  \item Produce a unified exclusion plot in the $(\Lambda, \alpha)$ plane
    combining all experimental channels.
\end{enumerate}
}

\outcome{Lower bound on $\Lambda$ from laboratory experiments; unified
exclusion plot; identification of the most sensitive experiment for
testing SCT at each distance scale.}

\deps{NT-1b, NT-4a (linearized equations for short-distance potential),
PPN-1 (Solar System).}

\textbf{Status (2026-03-17):} \textcolor{proven}{COMPLETE.}
Primary bound: $\Lambda > 2.565 \times 10^{-3}$~eV (95\% CL) from E\"ot-Wash torsion balance.
Corresponding Yukawa range $\lambda_1 < 35.8$~$\mu$m.
SCT potential is parameter-free: $\alpha_1 = -4/3$, $\alpha_2 = +1/3$, $V(0) = 0$.
No other experiment constrains SCT. 7.8\% tighter than PPN-1.
85 tests PASS, 5 figures. See \texttt{LT3d\_laboratory.tex}.

\subsubsection{LT-3e: Stellar Structure and TOV limits}
\label{task:LT3e}

\task{Derive the modified Tolman--Oppenheimer--Volkoff (TOV) equations for SCT and compute the maximum mass of neutron stars.}

\method{
\begin{enumerate}[(a)]
  \item Set up the interior static, spherically symmetric metric with a perfect fluid source.
  \item Treat the non-local higher-derivative terms perturbatively in $1/\Lambda^2$ (or use auxiliary scalar fields) to bypass the intractability of exact non-local integro-differential equations inside matter.
  \item Integrate the modified TOV equations for realistic nuclear equations of state (e.g., SLy, AP4).
  \item Verify that the non-local form factors do not artificially lower the maximum mass below the observational limit of $\sim 2.14 M_\odot$.
\end{enumerate}
}

\outcome{Maximum neutron star mass as a function of $\Lambda$; theoretical viability check against heavy pulsar observations.}

\deps{NT-4b (full nonlinear equations are required for strong-field interior metrics).}

\textbf{Severity: CRITICAL.} (Direct astrophysics falsification point).

% ==========================================================================
\section{Lessons from the Analysis of the Combined Approach (Part IV)}
\label{sec:lessons}
% ==========================================================================

The analysis of a combined approach (referred to as the ``stitched'' model,
analogous to an engineering assembly of foreign modules) in Part~IV yielded
both cautionary lessons and constructive insights for SCT.

\subsection{Structural critique of the stitched approach}

The combined model assembles three independent mechanisms:
\begin{enumerate}
  \item \textbf{Vacuum energy sequestration} (Kaloper--Padilla type):
    a global constraint decouples the cosmological constant from UV physics.
  \item \textbf{Ghost-free nonlocality} (Tomboulis--Modesto--Biswas--Mazumdar--Siegel type):
    entire-function form factors tame UV divergences.
  \item \textbf{Islands/QES} (Ryu--Takayanagi--Engelhardt--Wall type):
    the island formula for entanglement entropy.
\end{enumerate}

\textbf{Key deficiencies identified:}
\begin{itemize}
  \item \textbf{No single generating principle.} The three sectors are simply added,
    with no explanation of why these particular mechanisms coexist or are compatible.
  \item \textbf{UV-completion not specified.} The model explicitly marks
    UV-completion as an open question, offering no path forward.
  \item \textbf{Spectral dimension $d_S: 4 \to 2$ not mentioned.}
    A universal feature of quantum gravity approaches is absent.
  \item \textbf{Black hole entropy not derived.} The Bekenstein--Hawking formula
    is assumed, not obtained from the framework.
  \item \textbf{No new predictions beyond component models.} The combination
    does not generate predictions that its components individually lack.
\end{itemize}

\subsection{Strategic contrast: SCT vs.\ stitched approach}

\begin{center}
\begin{tabular}{lcc}
\toprule
\textbf{Feature} & \textbf{Stitched model} & \textbf{SCT} \\
\midrule
Generating principle & Three separate & One (spectral triple) \\
Form factors & Postulated & Must be derived \\
Sequestration & Added & Potentially automatic \\
BH entropy & Assumed & Must be derived \\
$d_S$ flow & Not addressed & Must be computed \\
UV-completion & ``Not specified'' & Must be proven \\
\bottomrule
\end{tabular}
\end{center}

\subsection{Constructive lessons for SCT}

Despite structural deficiencies, the stitched model identifies useful
directions:

\begin{enumerate}
  \item \textbf{Prediction format.} The three channels---smoothed Newtonian potential
    $\Delta F/F(r)$, QNM frequency shifts $\delta\omega$, and cosmological equation
    of state $w(z)$---are the correct prediction targets. SCT should produce numbers
    in these same channels.

  \item \textbf{Entire form factors as effective description.} The postulated
    ghost-free form factors work well as a phenomenological model. In SCT, the key
    question is whether these arise \emph{naturally} from the spectral triple
    (Task~NT-2).

  \item \textbf{Sequestration as a hint.} The spectral action
    $\Tr(f(D^2/\Lam^2))$ may be automatically insensitive to additive shifts
    in the vacuum energy, providing a natural sequestration mechanism without
    a separate postulate. This should be investigated.

  \item \textbf{The ``code'' model.} The relation
    $S(A) = \text{Area}(\gamma_A)/(4G_{\text{eff}}) + S_{\text{bulk}}$
    as a \emph{definition} of geometry from entanglement data is closest to the
    SCT spirit, where geometry is spectral data. Exploring this connection is
    a priority for Task~LT-2.
\end{enumerate}

% ==========================================================================
\section{Priority Order and Dependencies}
\label{sec:priority}
% ==========================================================================

\subsection{Dependency graph}

The tasks form the following dependency structure (arrows indicate ``must be
completed before''):

\bigskip
\noindent\textbf{Level 0 (Foundation --- proven results):}
\begin{quote}
  Spectral action $\to$ Einstein--Hilbert limit [Eq.~\eqref{eq:EH-limit}] \\
  Ratio $c_1/c_2 = -1/3 + 120(\xi-1/6)^2/13$ [Eq.~\eqref{eq:ratio}; $= -1/3$ at $\xi = 1/6$] \\
  Master function $f(\xi)$ [Eq.~\eqref{eq:master}]
\end{quote}

\noindent\textbf{Level 1 (Near-term + independent):}
\begin{quote}
  \textbf{NT-1} (Dirac form factors) $\longleftarrow$ Level 0
    \hfill\textcolor{proven}{[\,COMPLETE\,]} \\[2pt]
  \textbf{NT-1b} (Scalar + vector + SM form factors) $\longleftarrow$ NT-1
    \hfill\textcolor{proven}{[\,PHASES 1--3 COMPLETE\,]} \\[2pt]
  \textbf{NT-2} (Entire-function property) $\longleftarrow$ NT-1 \\[2pt]
  \textbf{NT-3} (Spectral dimension) $\longleftarrow$ Level 0 (partial); MT-3 (full) \\[2pt]
  \textbf{NT-4a} (Linearized nonlocal eqs.) $\longleftarrow$ NT-1, NT-1b \\[2pt]
  \textbf{META-1} (Parameter counting + $f$-analysis) $\longleftarrow$ Level 0 (iterative) \\[2pt]
  \textbf{MR-1} (Lorentzian spectral action) $\longleftarrow$ Level 0
    \quad\textcolor{long}{[CRITICAL]} \\[2pt]
  \textbf{MR-6} (Curvature expansion convergence) $\longleftarrow$ Level 0
\end{quote}

\noindent\textbf{Level 1.5 (NT-4 subtasks):}
\begin{quote}
  \textbf{NT-4b} (Full nonlinear eqs.) $\longleftarrow$ NT-4a, NT-1, NT-1b \\[2pt]
  \textbf{NT-4c} (FLRW reduction) $\longleftarrow$ NT-4b, NT-1b
\end{quote}

\noindent\textbf{Level 2 (Mid-term and consistency):}
\begin{quote}
  \textbf{MT-1} (BH entropy + all 4 laws) $\longleftarrow$ NT-1, NT-1b, NT-4b \\[2pt]
  \textbf{MT-2} (Cosmology + $H_0$/$S_8$) $\longleftarrow$ NT-1, NT-1b, NT-2, NT-4c \\[2pt]
  \textbf{MT-3} (Postulate 5 variants) $\longleftarrow$ NT-1, NT-1b, NT-2 \\[2pt]
  \textbf{INF-1} (Inflation: $n_s$, $r$, $f_{\text{NL}}$) $\longleftarrow$ NT-1b, NT-2, NT-4c \\[2pt]
  \textbf{PPN-1} (PPN parameters) $\longleftarrow$ NT-1b, NT-4a
    \quad\textcolor{long}{[CRITICAL]} \\[2pt]
  \textbf{FUND-1} (QFT on curved spacetime) $\longleftarrow$ NT-1b, NT-4b, MT-3 \\[2pt]
  \textbf{MR-2} (Unitarity + stability) $\longleftarrow$ NT-1, NT-1b, NT-2, NT-4a
    \quad\textcolor{proven}{[CLOSED]} \\[2pt]
  \textbf{MR-3} (Causality + macrocausality) $\longleftarrow$ NT-1, NT-1b, NT-2, NT-4a
    \quad\textcolor{long}{[CRITICAL]} \\[2pt]
  \textbf{MR-7} (Scattering amplitudes) $\longleftarrow$ NT-1b, NT-4a, MR-2, MR-3 \\[2pt]
  \textbf{MR-9} (BH singularity + Penrose) $\longleftarrow$ NT-1, NT-1b, NT-2, NT-4b, MT-1 \\[2pt]
  \textbf{COMP-1} (Cross-program comparison) $\longleftarrow$ (iterative, all prediction tasks)
\end{quote}

\noindent\textbf{Level 3 (Long-term):}
\begin{quote}
  \textbf{MR-4} (Two-loop) $\longleftarrow$ NT-1, NT-1b, MR-2, MR-7
    \quad\textcolor{proven}{[UNCONDITIONAL]} \\[2pt]
  \textbf{MR-5} (All-orders + naturalness) $\longleftarrow$ NT-1, NT-1b, NT-2, MR-4
    \quad\textcolor{near}{[CONDITIONAL BV-3,4]} \\[2pt]
  \textbf{MR-8} (Non-perturbative definition) $\longleftarrow$ MT-3, MR-1 \\[2pt]
  \textbf{LT-1} (UV-completeness) $\longleftarrow$ NT-1, NT-1b, NT-2, MT-3,
    MR-2, MR-4, MR-5 \\[2pt]
  \textbf{LT-2} (Information paradox + firewalls) $\longleftarrow$ MT-1, MT-3, LT-1 \\[2pt]
  \textbf{LT-3a} (QNM predictions) $\longleftarrow$ NT-1b, NT-4b, MT-1, MR-9 \\[2pt]
  \textbf{LT-3b} (GW propagation) $\longleftarrow$ NT-4a, NT-4c, INF-1 \\[2pt]
  \textbf{LT-3c} (Cosmological predictions) $\longleftarrow$ MT-2, NT-1b, NT-4c \\[2pt]
  \textbf{LT-3d} (Laboratory/Solar System) $\longleftarrow$ NT-1b, NT-4a, PPN-1
\end{quote}

\subsection{Critical path}

The critical path (longest chain of dependent tasks) for UV-completeness is:
\begin{equation*}
  \boxed{\text{NT-1}\checkmark} \;\longrightarrow\;
  \boxed{\text{NT-1b}} \;\longrightarrow\;
  \boxed{\text{NT-4a}} \;\longrightarrow\;
  \boxed{\text{MR-2}} \;\longrightarrow\;
  \boxed{\text{MR-7}} \;\longrightarrow\;
  \boxed{\text{MR-4}} \;\longrightarrow\;
  \boxed{\text{MR-5}} \;\longrightarrow\;
  \boxed{\text{LT-1}}
\end{equation*}
\noindent(NT-2, also required by MR-2, runs in parallel with NT-1b $\to$ NT-4a
and does not lengthen this path.)

A parallel critical path for the physical predictions track (cosmological):
\begin{equation*}
  \boxed{\text{NT-1}\checkmark} \;\longrightarrow\;
  \boxed{\text{NT-1b (Ph.\,1--2}\checkmark\text{)}} \;\longrightarrow\;
  \boxed{\text{NT-4a}} \;\longrightarrow\;
  \boxed{\text{NT-4b}} \;\longrightarrow\;
  \boxed{\text{NT-4c}} \;\longrightarrow\;
  \boxed{\text{INF-1}} \;\longrightarrow\;
  \boxed{\text{LT-3b}}
\end{equation*}

A third critical path for Solar System/laboratory tests (shortest):
\begin{equation*}
  \boxed{\text{NT-1}\checkmark} \;\longrightarrow\;
  \boxed{\text{NT-1b}} \;\longrightarrow\;
  \boxed{\text{NT-4a}} \;\longrightarrow\;
  \boxed{\text{PPN-1}} \;\longrightarrow\;
  \boxed{\text{LT-3d}}
\end{equation*}

\textbf{Key advantage of the NT-4 split:} The Solar System path now requires only
NT-4a (linearized equations), not the full NT-4b/c. This de-risks the most
experimentally urgent critical path.

NT-1 is \textcolor{proven}{COMPLETE}.
NT-1b Phases~1--2 (scalar and vector form factors) are
\textcolor{proven}{COMPLETE}; Phase~3 ($c_1/c_2$ re-derivation with full SM
spectrum) remains \textbf{pending}.
Once Phase~3 finishes, NT-1b unblocks NT-4a (linearized equations), MR-2 (propagator),
and MR-7 (amplitudes).
After NT-1b, \textbf{NT-4a (linearized nonlocal equations)} feeds directly into
PPN-1, MR-2, MR-3, and MR-7, while \textbf{NT-4b} (full nonlinear) feeds into
MT-1, MR-9, and \textbf{NT-4c} (FLRW reduction) feeds into MT-2, INF-1.
In parallel, \textbf{MR-1 (Lorentzian spectral action)} and
\textbf{MR-6 (convergence)} can proceed independently.

\subsection{Recommended execution order}

\begin{longtable}{clp{3.2cm}p{4.5cm}}
\toprule
\textbf{Pri.} & \textbf{Task} & \textbf{Blocked by} & \textbf{Unblocks} \\
\midrule
\endfirsthead
\toprule
\textbf{Pri.} & \textbf{Task} & \textbf{Blocked by} & \textbf{Unblocks} \\
\midrule
\endhead
--- & \textcolor{proven}{NT-1: Dirac form factors} & \textcolor{proven}{DONE} & \textcolor{proven}{(all below)} \\
\midrule
1 & META-1: Theory census + $f$-analysis & --- & (foundation for all reviews) \\
2 & \textcolor{proven}{NT-1b: Ph.\,1--3 done} & \textcolor{proven}{DONE} & NT-4a, MR-2, MR-4, MR-7, MT-1, MT-2, INF-1, PPN-1, LT-3a--e \\
3 & MR-1: Lorentzian action & --- & MR-8 \\
4 & \textcolor{proven}{MR-6: Convergence} & \textcolor{proven}{DONE} & (cross-check) \\
5 & NT-2: Entire functions & NT-1 & MR-2, MR-3, MT-2, MT-3, MR-5, INF-1 \\
6 & \textcolor{proven}{NT-3: Spectral dimension} & \textcolor{proven}{DONE} & (interpretation) \\
7a & NT-4a: Linearized nonlocal eqs. & NT-1, NT-1b & MR-2, MR-3, MR-7, PPN-1, NT-4b \\
7b & NT-4b: Full nonlinear eqs. & NT-4a & MT-1, MR-9, NT-4c, FUND-1, LT-3a \\
7c & NT-4c: FLRW reduction & NT-4b, NT-1b & MT-2, INF-1, LT-3b, LT-3c \\
8 & \textcolor{proven}{MR-2: Unitarity+stability [CLOSED]} & \textcolor{proven}{DONE} & MR-4, MR-7, LT-1 \\
9 & MR-3: Causality+macrocausality & NT-1b, NT-2, NT-4a & MR-7 \\
10 & PPN-1: Post-Newtonian limits & NT-4a, NT-1b & LT-3d (verification) \\
11 & \textcolor{proven}{MT-2: Cosmology+tensions [NEGATIVE]} & NT-1b, NT-2, NT-4c & LT-3b, LT-3c \\
12 & \textcolor{near}{MT-1: BH entropy+4 laws [CONDITIONAL]} & NT-1b, NT-4b & LT-2, LT-3a, MR-9 \\
13 & \textcolor{near}{INF-1: Spectral inflation [CONDITIONAL]} & NT-1b, NT-2, NT-4c & LT-3b, LT-3c \\
13b & \textcolor{long}{INF-2: Dilaton inflation [NEGATIVE]} & INF-1, NT-1b, NT-4b & (output) \\
14 & MT-3: Postulate 5 & NT-1, NT-1b, NT-2 & LT-1, LT-2, MR-8, FUND-1 \\
15 & MR-7: Scattering ampl. & NT-1b, NT-4a, MR-2, MR-3 & MR-4 \\
16 & MR-9: BH singularity+Penrose & NT-1b, NT-2, NT-4b, MT-1 & LT-3a \\
17 & FUND-1: Axiomatic foundations & NT-1b, NT-4b, MT-3 & (structural integrity) \\
18 & \textcolor{proven}{MR-4: Two-loop [UNCONDITIONAL]} & \textcolor{proven}{DONE} & MR-5, LT-1 \\
19 & \textcolor{near}{MR-5: All-orders [COND BV-3,4]} & NT-1, NT-1b, NT-2, MR-4 & LT-1 \\
20 & MR-8: Non-perturbative & MT-3, MR-1 & (structural integrity) \\
21 & LT-3a: QNM predictions & NT-1b, NT-4b, MT-1, MR-9 & (output) \\
22 & LT-3b: GW propagation & NT-4a, NT-4c, INF-1 & (output) \\
23 & LT-3c: Cosmological predictions & MT-2, NT-1b, NT-4c & (output) \\
24 & LT-3d: Laboratory/Solar System & NT-1b, NT-4a, PPN-1 & (output) \\
25 & LT-3e: Stellar structure (TOV) & NT-4b & (output) \\
26 & LT-1: UV-completeness & NT-1b, NT-2, MT-3, MR-2, MR-4, MR-5 & LT-2 \\
27 & LT-2: Info.\ paradox+firewalls & MT-1, MT-3, LT-1 & (final output) \\
\midrule
\multicolumn{4}{l}{\footnotesize COMP-1 (cross-program comparison) runs iteratively throughout, updated as results accumulate.} \\
\bottomrule
\end{longtable}

\subsection{Parallelization opportunities}

The following tasks can proceed simultaneously:
\begin{itemize}
  \item \textbf{Now (NT-1 complete):} NT-1b, NT-2, MR-1, MR-6, NT-3,
    \textbf{META-1} are all independent of each other.
    META-1 (theory census + cutoff function analysis) has no dependencies
    and should start immediately.
  \item \textbf{After NT-1b:} NT-4a (linearized equations) can start
    (needs only NT-1 + NT-1b).
  \item \textbf{After NT-4a:} \textbf{PPN-1} (post-Newtonian limits) can start
    immediately, in parallel with MR-2, MR-3, and NT-4b.
    PPN-1 is \textbf{CRITICAL}---a Solar System test of the linearized potential.
    \emph{Key advantage of the NT-4 split:} PPN-1 no longer waits for the full
    nonlinear variational problem (NT-4b).
  \item \textbf{After NT-4a + NT-2:} MR-2 and MR-3 can proceed in parallel.
    MR-7 follows once MR-2 + MR-3 are both complete (depends on both).
  \item \textbf{After NT-4b:} NT-4c (FLRW reduction) starts, leading to MT-2 and INF-1.
  \item \textbf{After NT-4c:} MT-2 and INF-1 can proceed in parallel
    (INF-1 no longer depends on MT-2).
  \item \textbf{MT-1 and MT-2} are independent (MT-1 needs NT-4b; MT-2 needs NT-4c).
  \item \textbf{MR-7 $\to$ MR-4:} MR-7 (scattering amplitudes) must now precede
    MR-4 (two-loop), as the one-loop integrand structure feeds the two-loop calculation.
  \item \textbf{LT-3a/b/c/d} are largely independent of each other:
    LT-3a (QNMs) needs MT-1; LT-3b (GW propagation) needs INF-1;
    LT-3c (cosmology) needs MT-2; LT-3d (lab tests) needs PPN-1.
  \item \textbf{LT-1 and LT-3a--d} are independent (LT-3 does not require LT-1).
  \item \textbf{MR-8 and MR-4} are independent (different dependency chains).
  \item \textbf{MR-9} can proceed as soon as NT-4b + MT-1 are done,
    independently of the UV-completeness track.
  \item \textbf{FUND-1} can proceed after NT-4b + MT-3, independently of prediction tasks.
  \item \textbf{COMP-1} runs iteratively---updated after each major result,
    no blocking dependencies.
\end{itemize}

\subsection{Methods and tools per task}

\begin{longtable}{clp{6.5cm}}
\toprule
\textbf{Task} & \textbf{Primary tools} & \textbf{Verification} \\
\midrule
\endfirsthead
\toprule
\textbf{Task} & \textbf{Primary tools} & \textbf{Verification} \\
\midrule
\endhead
NT-1 & SymPy, SciPy, \texttt{sct\_tools} & \textcolor{proven}{DONE.} 274+ checks PASS (7 audit rounds). \\
NT-1b & SymPy, SciPy, \texttt{sct\_tools}, cadabra2\textsuperscript{W} & Local limits vs.\ known $\beta$-functions; consistency with NT-1 Dirac results \\
NT-2 & Complex analysis (SymPy), mpmath & Hadamard factorization; numerical pole search \\
NT-3 & \texttt{sct\_tools.graphs}, networkx, SciPy & Check $d_S \to 4$ in IR; compare with CDT lattice results \\
NT-4a & SymPy (linearized), \texttt{sct\_tools} & Flat-space limit $\to$ massless graviton; gauge invariance check \\
NT-4b & cadabra2\textsuperscript{W} (variational) & Local limit $\to$ Einstein eqs.; Bianchi identity; trace consistency \\
NT-4c & SymPy, SciPy (FRW reduction) & de Sitter limit; GR Friedmann recovery; perturbation stability \\
\midrule
META-1 & Manual analysis, LaTeX tables & 10-check review; competitor comparison; falsifiability matrix \\
\midrule
MR-1 & Formal (Krein space theory) & Wick rotation consistency; compare Paschke--Sitarz, van den Dungen \\
MR-2 & SymPy, cadabra2\textsuperscript{W} (propagator) & Residue signs; optical theorem; comparison with Stelle, Tomboulis \\
MR-3 & SymPy + SciPy (Green's fn.) & Kramers--Kronig; retarded propagator; macrocausality (Krasnikov test) \\
MR-4 & cadabra2\textsuperscript{W}, pySecDec\textsuperscript{W} (2-loop) & Compare Goroff--Sagnotti; cross-check divergence cancellation \\
MR-5 & Formal (power-counting) & Explicit check at 2-loop (MR-4); compare Tomboulis, Modesto arguments \\
MR-6 & SymPy + numerical & Exact $S^4$ spectral action vs.\ truncated expansion \\
MR-7 & cadabra2\textsuperscript{W}, SymPy (amplitudes) & Ward identities; IR limit = GR; unitarity cuts \\
MR-8 & Formal + numerical (matrix models) & Barrett--Glaser comparison; CDT cross-check; Poisson sprinkling convergence (Glaser--Stern) \\
MR-9 & SciPy (ODE integration) & Kretschmer scalar finiteness; Penrose theorem; Raychaudhuri defocusing \\
\midrule
PPN-1 & SymPy (linearized eqs.), SciPy & Cassini bound; LLR; Mercury perihelion; binary pulsars; E\"ot--Wash \\
INF-1 & SymPy, CAMB, CLASS\textsuperscript{W}, cobaya & Planck $n_s$/$r$; BICEP/Keck; LiteBIRD forecasts; non-Gaussianity \\
FUND-1 & Formal analysis, SymPy & Born--Oppenheimer sep.; Haag--Kastler axioms; anomaly check \\
COMP-1 & Literature analysis, LaTeX tables & Feature-by-feature comparison; domain map vs.\ 5 programs \\
\midrule
MT-1 & SymPy (Wald entropy), healpy\textsuperscript{W} & Compare with Bekenstein--Hawking; LQG $c_{\log}$; 4 laws \\
MT-2 & CAMB, CLASS\textsuperscript{W}, healpy\textsuperscript{W}, emcee & Fit to Planck + DESI + Pantheon+; $H_0$ and $S_8$ tensions; $\Lambda_{\text{CC}}$ consistency (spectral vs.\ causal set) \\
MT-3 & Formal analysis, SymPy & Well-posedness proofs; classical limit check \\
\midrule
LT-1 & Functional RG (numerical), pySecDec\textsuperscript{W} & Fixed-point search; requires MR-2, MR-4, MR-5 inputs \\
LT-2 & Formal + numerical & Page curve; firewall resolution (Giddings-type); unitarity check \\
LT-3a & SciPy (ODE), cadabra2\textsuperscript{W} & Modified Regge--Wheeler/Zerilli; Kerr via Teukolsky; LIGO/ET data \\
LT-3b & SymPy, SciPy (dispersion) & GW speed; polarization; secondary spectrum; LISA forecasts \\
LT-3c & CLASS\textsuperscript{W}, cobaya, healpy\textsuperscript{W} (MCMC) & $w(z)$ from hi\_class; DESI + Euclid comparison \\
LT-3d & SciPy (numerical), SymPy & E\"ot--Wash torsion balance; Casimir; unified exclusion plot \\
LT-3e & SciPy (ODE integration) & Max neutron star mass $\geq 2.14 M_\odot$ \\
\bottomrule
\end{longtable}

\noindent\textsuperscript{W}~Runs in WSL2 Ubuntu~24.04 (venv \texttt{\textasciitilde/sct-wsl}).
cadabra2: tensor CAS (index contractions, Bianchi identities, variational calculus).
pySecDec: multi-loop Feynman integrals via sector decomposition.
healpy: HEALPix CMB analysis (angular power spectra, $C_\ell$ estimation).
CLASS: Boltzmann solver (alternative to CAMB for cross-validation).

% ==========================================================================
\section*{Summary}
% ==========================================================================

\textbf{Completed:} NT-1 (Dirac form factors $F_1$, $F_2$) and
NT-1b Phases~1--3 (scalar + vector + combined SM form factors) are complete.
Total: $4156{+}$ quantitative checks passing, 71/71~PDFs compiled,
41~Lean identities with 46~proof files (91~Lean tests).
Three papers published (DOIs: 10.5281/zenodo.19039242,
10.5281/zenodo.19045796, 10.5281/zenodo.19056204).

\textbf{Current status:} The roadmap contains \textbf{31 tasks}
(27 original, with NT-4 split adding 2, LT-3e adding 1; plus COMP-1 running iteratively)
organized in five categories:
\begin{itemize}
  \item \textbf{Near-term (5+2 subtasks):} NT-1 \textcolor{proven}{[\,DONE\,]},
    NT-1b \textcolor{proven}{[\,Ph.\,1--3 DONE\,]},
    NT-2 \textcolor{proven}{[\,DONE\,]}, NT-3 \textcolor{proven}{[\,DONE\,]},
    NT-4a \textcolor{proven}{[\,DONE\,]}, NT-4b \textcolor{proven}{[\,DONE\,]},
    NT-4c \textcolor{proven}{[\,DONE\,]}.
  \item \textbf{Mid-term (3+2):} MT-1 (BH entropy + 4 laws), MT-2 (cosmology + tensions),
    MT-3 (Postulate~5); plus \textbf{INF-1} (spectral inflation) and
    \textbf{META-1} (theory census + cutoff function $f$ analysis).
  \item \textbf{Theory consistency (9+2):} MR-1 through MR-9; plus
    \textbf{PPN-1} (post-Newtonian limits, \textbf{CRITICAL}) and
    \textbf{FUND-1} (axiomatic foundations + fermion doubling).
    MR-1, MR-2, MR-3, MR-4, MR-5, and PPN-1 are \textbf{CRITICAL} for theory viability.
  \item \textbf{Long-term (7):} LT-1 (UV-completeness), LT-2 (information
    paradox + firewalls), LT-3a (QNMs), LT-3b (GW propagation),
    LT-3c (cosmological predictions), LT-3d (laboratory/Solar System tests),
    LT-3e (stellar structure/TOV limits).
  \item \textbf{Iterative:} COMP-1 (cross-program comparison + quantitative AS
    comparison, updated continuously).
\end{itemize}

\textbf{Immediate priorities (remaining open tasks):}
\begin{enumerate}
  \item \textbf{META-1}: Theory census, falsifiability matrix, and cutoff
    function $f$ sensitivity analysis (no dependencies; foundational self-audit).
  \item \textbf{MT-3}: Three variants of Postulate~5 (quantum formulation).
  \item \textbf{MR-8}: Non-perturbative definition (path integral over Dirac operators).
  \item \textbf{MR-9}: Black hole singularity resolution.
  \item \textbf{FUND-1}: Recovery of QFT on curved spacetime.
  \item \textbf{LT-1}: UV-completeness (requires MT-3, MR-5).
  \item \textbf{LT-3a--e}: Testable numerical predictions (QNMs, GWs, cosmology, lab, TOV).
  \item \textbf{LT-2}: Information paradox and island formula (requires MT-1, MT-3, LT-1).
  \item \textbf{COMP-1}: Cross-program comparison (iterative, updated continuously).
\end{enumerate}

\textbf{Completed CRITICAL tasks:}
\begin{itemize}
  \item \textbf{MR-1} \textcolor{proven}{[\,COMPLETE\,]}: Lorentzian spectral action.
  \item \textbf{MR-2} \textcolor{proven}{[\,CLOSED\,]}: Unitarity ($D^2$-quantization, bounded propagator, no ghosts).
  \item \textbf{MR-3} \textcolor{near}{[\,CONDITIONAL\,]}: Causality (macrocausal; microcausality violated at $\sim 0.884/\Lambda$).
  \item \textbf{MR-4} \textcolor{proven}{[\,UNCONDITIONAL\,]}: Two-loop (CHIRAL-Q chirality theorem).
  \item \textbf{MR-5} \textcolor{near}{[\,CONDITIONAL BV-3,4\,]}: Finiteness (UV-finite in $D^2$-quantization).
  \item \textbf{PPN-1} \textcolor{proven}{[\,COMPLETE\,]}: Solar System tests ($\Lambda \geq 2.38 \times 10^{-3}$\,eV).
\end{itemize}

\textbf{Three critical paths:}
\begin{enumerate}
  \item \textbf{UV-completeness:} NT-1$\checkmark$ $\to$ NT-1b $\to$ NT-4a $\to$
    MR-2 $\to$ MR-7 $\to$ MR-4 $\to$ MR-5 $\to$ LT-1 (NT-2 parallel with NT-1b$\to$NT-4a)
  \item \textbf{Cosmological predictions:} NT-1$\checkmark$ $\to$ NT-1b $\to$
    NT-4a $\to$ NT-4b $\to$ NT-4c $\to$ INF-1 $\to$ LT-3b
    (LT-3c shares NT-4c but reaches it via MT-2, not INF-1)
  \item \textbf{Solar System tests (shortest):} NT-1$\checkmark$ $\to$ NT-1b $\to$
    NT-4a $\to$ PPN-1 $\to$ LT-3d
\end{enumerate}

\textbf{Key datasets:}
\begin{itemize}
  \item CMB: Planck 2018, ACT DR6, Simons Observatory (upcoming), LiteBIRD (upcoming)
  \item BAO: DESI DR1/DR2, Euclid (upcoming)
  \item SNe: Pantheon+, Union3
  \item GW: LVK O4/O5 (ringdown, speed), LISA (upcoming), Einstein Telescope (upcoming)
  \item Dispersion: Fermi GBM, IceCube/KM3NeT (neutrino timing)
  \item Weak lensing: KiDS-1000, DES-Y3, HSC
  \item CMB $r$: BICEP/Keck, Simons Observatory, LiteBIRD
  \item PPN: Cassini, LLR, Mercury perihelion, binary pulsars, E\"ot-Wash
\end{itemize}

The strategic advantage of SCT over combined/stitched approaches lies in its
\emph{single generating principle}: all structure---form factors, ghost-freedom,
entropy, dimensional flow---must emerge from one object, the spectral triple.
The 11 theory-consistency tasks (MR-1 through MR-9, PPN-1, FUND-1) are
specifically designed to test whether this claim survives the requirements
for a genuine UV-complete theory.
Any single CRITICAL task failure would require fundamental revision of the theory.

\end{document}
